%=====================================================================
%  EXHAUSTIVE (line-by-line justified) version of the bounded
%  reduction section.
%
%  This file is a section to be \input into a master document whose
%  preamble defines:
%     \R \Om \Otilde \Fr \Hn \Per \eps \abs \norm \diam \dx
%  and the theorem-like environments
%     theorem/proposition/lemma/corollary/definition/remark
%  numbered by section, together with the bibliography keys BDV, FMP.
%  The cross-reference \ref{thm:boundedSV} resolves in the master file.
%
%  MATHEMATICAL CONTENT IS UNCHANGED relative to the verified source
%  sec_bounded_reduction_body.tex: same theorems, lemmas, labels,
%  logical order, and constants.  This file only ANNOTATES and EXPANDS;
%  it introduces no new claims and changes no constant.
%=====================================================================

\section{Bounded reduction}\label{sec:bddred}

The purpose of this section is to remove the boundedness hypothesis from
the sharp Saint--Venant stability inequality. The idea is simple to
state. We already possess a sharp quadratic stability estimate for sets
\emph{confined to a ball of a fixed radius}; we want the same estimate
for arbitrary open sets of finite measure, with a constant depending
only on the dimension. The obstruction is that a general competitor may
have very long, thin tentacles reaching out to infinity. The remedy,
following Brasco--De~Philippis--Velichkov, is to amputate those tails:
we will show that a set with small deficit carries almost all of its
torsional energy and almost all of its mass inside a ball of
dimensionally bounded radius, so that chopping it off there changes the
deficit and the asymmetry by only a controlled amount. The bounded
estimate, applied to the amputated set, then transfers back.

The input is the following \emph{bounded} estimate, proved in the earlier
part of the manuscript (Theorem~\ref{thm:boundedSV}), restated here in
energy form.

\begin{theorem}[Bounded sharp Saint--Venant]\label{thm:SV-bounded}
For every $n\ge2$ and every $R\ge2$ there is a constant
$c_{\mathrm{SV}}(n,R)>0$ such that
\begin{equation}\label{eq:SV-bounded}
   E(\Om)-E(B_1)\;\ge\;c_{\mathrm{SV}}(n,R)\,\Fr(\Om)^2
\end{equation}
for every open set $\Om\subset B_R$ with $\abs{\Om}=\omega_n$.
\end{theorem}

This is Theorem~\ref{thm:boundedSV}, restated in energy form: with the
constant $C(n,R,\rstar)$ there one may take
$c_{\mathrm{SV}}(n,R)=1/C(n,R,\rstar)$ and
$\dT(\Om)=E(\Om)-E(B_1)$, and the parameter $\rstar$ is fixed at, say,
$\rstar=\tfrac34$.
%% JUSTIFICATION.  The earlier theorem is phrased with the torsional
%% deficit \dT on the left and the asymmetry on the right separated by a
%% multiplicative constant C; writing the inequality with the constant on
%% the asymmetry side simply means inverting C.  The auxiliary parameter
%% \rstar is an internal parameter of that proof (an inner radius in the
%% selection step); any fixed admissible value works, and \rstar=3/4 is
%% one such.  No new content is introduced here.

Our goal is the global statement.

\begin{theorem}[Global sharp Saint--Venant stability]\label{thm:SV-global}
There is a dimensional constant $c_{\mathrm{SV}}^{\mathrm{glob}}(n)>0$
such that
\begin{equation}\label{eq:SV-global}
   E(\Om)-E(B_1)\;\ge\;c_{\mathrm{SV}}^{\mathrm{glob}}(n)\,\Fr(\Om)^2
\end{equation}
for every open set $\Om\subset\R^n$ with $\abs{\Om}=\omega_n$.
Equivalently, after rescaling, for every open set $\Om$ of finite
measure, with $B^*$ the ball of the same volume,
\begin{equation}\label{eq:T-global}
   T(B^*)-T(\Om)\;=\;2\bigl(E(\Om)-E(B^*)\bigr)\;\ge\;
   2\,c_{\mathrm{SV}}^{\mathrm{glob}}(n)\,
   \Bigl(\tfrac{\abs{\Om}}{\omega_n}\Bigr)^{(n+2)/n}\,\Fr(\Om)^2 .
\end{equation}
\end{theorem}

The strategy is the constructive reduction of Brasco, De~Philippis and
Velichkov \cite[\S5]{BDV}. From a set $\Om$ with small Saint--Venant
deficit we manufacture a \emph{bounded} surrogate $\Otilde$, of the same
volume, of dimensionally bounded diameter, with comparable deficit and
with controlled loss of asymmetry; applying
Theorem~\ref{thm:SV-bounded} to $\Otilde$ at a fixed dimensional radius
and transferring the estimate back to $\Om$ closes the small-deficit
regime. The complementary large-deficit regime is trivial because
$\Fr<2$. The whole construction is explicit: no compactness, no
minimising sequence and no contradiction argument enter.

The two technical ingredients are an $L^\infty$ tail estimate for the
torsion function (Lemma~\ref{lem:Linfty}, proved by a De~Giorgi
iteration) and the truncation construction itself
(Lemma~\ref{lem:truncation}). Both are proved here in full rather than
cited. The reader who wishes to keep the logical skeleton in mind may
note the order of dependence: a purely numerical convergence lemma
(\S\ref{ssec:DG-lemma}) feeds the De~Giorgi $L^\infty$ bound
(\S\ref{ssec:Linfty}), which feeds the energy-comparison inequality that
drives the truncation (\S\ref{ssec:truncation}); the truncation, the
trivial far branch (\S\ref{ssec:far}) and the transfer
(\S\ref{ssec:near}) are then assembled in \S\ref{ssec:global-proof}.

\subsection{Standing notation and the standard tools we cite}
\label{ssec:notation}

Throughout, $n\ge2$ and all sets are open subsets of $\R^n$ of finite
Lebesgue measure. We write $u_\Om\in H^1_0(\Om)$ for the torsion
function, i.e. the unique weak solution of $-\Delta u_\Om=1$ in $\Om$,
$u_\Om=0$ on $\partial\Om$, equivalently the minimiser of
\begin{equation}\label{eq:energy-def}
   J_\Om(v)\;=\;\tfrac12\int_\Om\abs{\nabla v}^2\dx-\int_\Om v\dx,
   \qquad v\in H^1_0(\Om),
\end{equation}
and $E(\Om)=J_\Om(u_\Om)=\min_{H^1_0(\Om)}J_\Om$.
%% JUSTIFICATION.  Existence and uniqueness of the minimiser are the
%% standard direct-method facts for the strictly convex, coercive
%% functional J_\Om on the Hilbert space H^1_0(\Om); the Euler--Lagrange
%% equation of J_\Om is exactly -\Delta u = 1 in the weak sense
%% \int \nabla u\cdot\nabla v = \int v for all v in H^1_0(\Om).
Testing the equation with $u_\Om$ gives the standard identities
\begin{equation}\label{eq:energy-identities}
   \int_\Om\abs{\nabla u_\Om}^2\dx=\int_\Om u_\Om\dx=T(\Om),
   \qquad
   E(\Om)=-\tfrac12\int_\Om u_\Om\dx=-\tfrac12\,T(\Om).
\end{equation}
%% JUSTIFICATION.  Put v=u_\Om in the weak equation: the left side is
%% \int|\nabla u_\Om|^2 and the right side is \int u_\Om; this is the
%% first identity, and both equal T(\Om) by definition of torsional
%% rigidity.  For the energy value, substitute v=u_\Om into J_\Om:
%%   E(\Om)=\tfrac12\int|\nabla u_\Om|^2-\int u_\Om
%%         =\tfrac12\int u_\Om-\int u_\Om=-\tfrac12\int u_\Om,
%% where the first \int u_\Om came from the just-proved identity.  Hence
%% E=-T/2.  In particular E(\Om)<0 for any nonempty \Om (the torsion
%% function is strictly positive on a nonempty open set by the maximum
%% principle, so \int u_\Om>0); we will use E(B_1)<0 repeatedly.

We use $\omega_n=\abs{B_1}$, $T_n=T(B_1)=\omega_n/(n(n+2))$, the Fraenkel
asymmetry $\Fr(\Om)=\inf_{x_0}\abs{\Om\,\triangle\,(B_1+x_0)}/\omega_n$
(for $\abs{\Om}=\omega_n$ this lies in $[0,2)$), and the scale-invariant
deficit
\begin{equation}\label{eq:deficit-def}
   D(\Om)\;=\;E(\Om)\abs{\Om}^{-(n+2)/n}-E(B_1)\abs{B_1}^{-(n+2)/n},
\end{equation}
so that $D(\Om)=\omega_n^{-(n+2)/n}\bigl(E(\Om)-E(B_1)\bigr)\ge0$ when
$\abs{\Om}=\omega_n$.
%% JUSTIFICATION (asymmetry range).  Trivially \Fr\ge0.  For the upper
%% bound, with B:=B_1+x_0 any unit ball and \abs{\Om}=\abs{B}=\omega_n,
%%   |\Om\triangle B| = |\Om\setminus B|+|B\setminus\Om|
%%                    = 2|\Om\setminus B| \le 2\omega_n,
%% using |\Om\setminus B|=|B\setminus\Om| (equal volumes) and
%% |\Om\setminus B|\le|\Om|=\omega_n; equality |\Om\triangle B|=2\omega_n
%% needs disjoint sets, excluded for an infimum over translates of a
%% fixed-volume set, so \Fr<2.
%% JUSTIFICATION (deficit value at \abs{\Om}=\omega_n).  Put
%% \abs{B_1}=\omega_n in \eqref{eq:deficit-def}: the second term is
%% E(B_1)\omega_n^{-(n+2)/n}, and the first is
%% E(\Om)\omega_n^{-(n+2)/n}, so D(\Om)=\omega_n^{-(n+2)/n}(E(\Om)-E(B_1)).
The non-negativity of $D$ is the
\emph{Saint--Venant inequality} (equivalently Talenti's comparison
theorem): among sets of given volume the ball minimises $E$, i.e.
\begin{equation}\label{eq:saint-venant}
   E(\Om)\,\abs{\Om}^{-(n+2)/n}\;\ge\;E(B)\,\abs{B}^{-(n+2)/n}
   \qquad\text{for every ball }B.
\end{equation}
We treat \eqref{eq:saint-venant} as a known classical result.
%% JUSTIFICATION.  The quantity E(\Om)\abs{\Om}^{-(n+2)/n} is invariant
%% under dilations because E(t\Om)=t^{n+2}E(\Om) (see the scaling note
%% below) and \abs{t\Om}=t^n\abs{\Om}; thus \eqref{eq:saint-venant} for
%% all balls B is equivalent to the single statement that B minimises
%% this scale-invariant energy among fixed-volume sets, which is the
%% Saint--Venant--Talenti theorem.

We will also need the elementary scaling law for $E$. If $\Om'$ is open
and $t>0$ then $E(t\Om')=t^{n+2}E(\Om')$, with torsion function
$u_{t\Om'}(x)=t^2u_{\Om'}(x/t)$.
%% JUSTIFICATION.  If -\Delta u_{\Om'}=1 in \Om' then setting
%% w(x):=t^2 u_{\Om'}(x/t) gives \Delta w(x)=t^2\cdot t^{-2}(\Delta
%% u_{\Om'})(x/t)=t^2\cdot t^{-2}\cdot(-1)=-1 in t\Om', and w=0 on
%% \partial(t\Om'); by uniqueness w=u_{t\Om'}.  Then
%%   E(t\Om')=-\tfrac12\int_{t\Om'}u_{t\Om'}(x)\,dx
%%           =-\tfrac12\int t^2 u_{\Om'}(x/t)\,dx
%%           =-\tfrac12\, t^2\cdot t^n\int_{\Om'}u_{\Om'}(y)\,dy
%%           =t^{n+2}E(\Om'),
%% changing variables y=x/t (Jacobian t^n).  We use this in Steps 6 and
%% in the general-volume rescaling.

The external analytic facts we invoke without proof are:
\begin{itemize}
\item the \emph{Sobolev inequality}
   $\norm{w}_{L^{2^*}(\R^n)}\le S_n\norm{\nabla w}_{L^2(\R^n)}$ for
   $w\in H^1(\R^n)$ when $n\ge3$ (with $2^*=2n/(n-2)$), and its $n=2$
   surrogate explained in Remark~\ref{rem:n2};
\item the \emph{Saint--Venant/Talenti inequality}~\eqref{eq:saint-venant},
   used in the truncation argument and in the global $L^\infty$ bound
   \eqref{eq:global-Linfty} below;
\item Theorem~\ref{thm:SV-bounded}, the bounded stability input;
\item the \emph{suboptimal} torsional stability bound
   $\Fr(\Om)^4\le C(n)D(\Om)$ (Lemma~\ref{lem:seed}), used only to seed
   the truncation iteration. This is proved here from the level-set identity
   established earlier in the manuscript; see Remark~\ref{rem:seed-status}.
\end{itemize}
Everything else --- in particular the fast-convergence iteration lemma
(Lemma~\ref{lem:DG-iteration2}), the De~Giorgi $L^\infty$ tail estimate
(Lemma~\ref{lem:Linfty}), the energy-comparison inequality, the truncation
iteration (Lemma~\ref{lem:truncation}) and the gluing --- is proved from
scratch.

We will also need the elementary pointwise sup-bound on the torsion
function. By Talenti's symmetrisation comparison the torsion function of
$\Om$ is dominated, after Schwarz symmetrisation, by that of the ball of
the same volume; in particular
\begin{equation}\label{eq:global-Linfty}
   \norm{u_\Om}_{L^\infty(\Om)}\;\le\;\norm{u_{B^*}}_{L^\infty(B^*)}
   \;=\;\frac{(\abs{\Om}/\omega_n)^{2/n}}{2n}
   \;=:\;\Lambda_n\,\abs{\Om}^{2/n},
   \qquad \Lambda_n:=\frac{\omega_n^{-2/n}}{2n},
\end{equation}
where $B^*$ is the ball with $\abs{B^*}=\abs{\Om}$ and we used
$u_{B_r}(x)=(r^2-\abs{x}^2)/(2n)$. For $\abs{\Om}=\omega_n$ this reads
$\norm{u_\Om}_{L^\infty(\Om)}\le 1/(2n)$. We record \eqref{eq:global-Linfty}
as a consequence of the cited Talenti comparison.
%% JUSTIFICATION (the ball's torsion function).  For \Om=B_r the function
%% u(x)=(r^2-|x|^2)/(2n) satisfies -\Delta u = -\tfrac1{2n}\Delta(r^2-|x|^2)
%% = -\tfrac1{2n}(-2n)=1 (since \Delta|x|^2=2n) and vanishes on |x|=r; by
%% uniqueness it is u_{B_r}.  Its maximum is at the centre, u(0)=r^2/(2n).
%% JUSTIFICATION (the sup-bound).  Talenti's pointwise comparison says
%% the decreasing rearrangement u_\Om^* is pointwise \le u_{B^*}, where
%% B^* is the ball of volume \abs{\Om}; in particular the maxima compare,
%% \|u_\Om\|_\infty = \|u_\Om^*\|_\infty \le \|u_{B^*}\|_\infty.  Writing
%% \abs{B^*}=\abs{\Om}=\omega_n(R^*)^n gives R^*=(\abs{\Om}/\omega_n)^{1/n}
%% and \|u_{B^*}\|_\infty=(R^*)^2/(2n)=(\abs{\Om}/\omega_n)^{2/n}/(2n),
%% which is \eqref{eq:global-Linfty}.  At \abs{\Om}=\omega_n the right
%% side is 1/(2n).  We will use the cruder bound u_\Om\le 1/(2n)\le1 in
%% the De~Giorgi iteration.

\subsection{The fast-convergence iteration lemma}\label{ssec:DG-lemma}

The De~Giorgi scheme reduces an $L^\infty$ bound to the vanishing of a
nonlinear recursion. We isolate the underlying numerical lemma and prove
it, so that nothing is cited at this step. The mechanism is purely
algebraic: a super-quadratic recursion (the increment $a_{k+1}$ depends
on a power $a_k^{1+\beta}$ strictly larger than the first power)
ultimately beats any \emph{fixed} geometric growth factor $b^k$, provided
the seed $a_0$ is below an explicit threshold. The threshold is exactly
the break-even point of the two competing effects.

\begin{lemma}[Fast geometric convergence]\label{lem:DG-iteration2}
Let $(a_k)_{k\ge0}$ be non-negative reals with
\begin{equation}\label{eq:DG-recursion2}
   a_{k+1}\;\le\;C\,b^{\,k}\,a_k^{\,1+\beta},\qquad k\ge0,
\end{equation}
where $C>0$, $b>1$, $\beta>0$. If
\begin{equation}\label{eq:DG-smallness2}
   a_0\;\le\;C^{-1/\beta}\,b^{-1/\beta^2},
\end{equation}
then
\begin{equation}\label{eq:DG-conclusion}
   a_k\;\le\;b^{-k/\beta}\,a_0\qquad\text{for all }k\ge0,
\end{equation}
so $a_k\to0$.
\end{lemma}

\begin{proof}
Put $r:=b^{-1/\beta}\in(0,1)$ and $\theta:=C\,a_0^{\beta}$.
%% JUSTIFICATION.  b>1 and 1/\beta>0 give b^{-1/\beta}\in(0,1), so r\in(0,1);
%% \theta\ge0 since C>0, a_0\ge0.
The smallness hypothesis \eqref{eq:DG-smallness2} is precisely
$\theta\le r$.
%% JUSTIFICATION.  Raise \eqref{eq:DG-smallness2} to the power \beta>0
%% (monotone for nonnegative bases):
%%   a_0^{\beta}\le C^{-1}b^{-1/\beta}=C^{-1}r,
%% i.e.\ C a_0^{\beta}\le r, which is \theta\le r.  (Equivalence: the
%% steps are reversible since x\mapsto x^{1/\beta} is increasing.)
We prove the target $a_k\le r^{k}a_0$ by induction; this is
\eqref{eq:DG-conclusion} because $r^k=b^{-k/\beta}$, and $r<1$ then gives
$a_k\to0$.

\emph{Base case} $k=0$: the claim reads $a_0\le r^0a_0=a_0$, an equality.

\emph{Inductive step.} Assume $a_k\le r^{k}a_0$. Using
\eqref{eq:DG-recursion2} and splitting off one factor $a_k$,
\[
   a_{k+1}\le C\,b^{k}\,a_k^{1+\beta}
          =C\,b^{k}\,a_k\,a_k^{\beta}
          \le C\,b^{k}\,a_k\,\bigl(r^{k}a_0\bigr)^{\beta}
          =\bigl(C\,a_0^{\beta}\bigr)\,(b\,r^{\beta})^{k}\,a_k
          =\theta\,(b\,r^{\beta})^{k}\,a_k .
\]
%% JUSTIFICATION, term by term.
%%  (i) a_k^{1+\beta}=a_k\cdot a_k^{\beta}: laws of exponents.
%%  (ii) a_k^{\beta}\le(r^k a_0)^{\beta}: the inductive hypothesis
%%       a_k\le r^k a_0 raised to the power \beta>0 (monotone).  Here we
%%       use a_k\ge0 and r^k a_0\ge0.
%%  (iii) regroup: C\,b^k\,(r^k a_0)^{\beta}
%%        =(C a_0^{\beta})\,b^k r^{k\beta}=(C a_0^{\beta})(b\,r^{\beta})^k.
%%  (iv) rename C a_0^{\beta}=\theta.
By the definition of $r$ we have $r^{\beta}=b^{-1}$, hence $b\,r^{\beta}=1$
and the geometric factor collapses:
\[
   a_{k+1}\le\theta\,a_k\le r\,a_k\le r\cdot r^{k}a_0=r^{k+1}a_0,
\]
using $\theta\le r$ and the inductive hypothesis $a_k\le r^k a_0$.
%% JUSTIFICATION.  r=b^{-1/\beta} gives r^{\beta}=b^{-1}, so b r^{\beta}=1
%% and (b r^{\beta})^k=1; this kills the only k-dependent factor.  Then
%% \theta\le r (smallness) and a_k\le r^k a_0 (induction) chain the
%% inequalities.  This is exactly the target at index k+1.
This closes the induction.
\end{proof}

\begin{remark}
The identity $b\,r^{\beta}=1$ with $r=b^{-1/\beta}$ is the algebraic
mechanism behind the convergence: the super-quadratic gain
$a_k^{1+\beta}$ converts, once $a_0$ is small, the growing factor $b^k$
into a contraction. The threshold \eqref{eq:DG-smallness2} is sharp for
this clean conclusion: it is exactly the value of $a_0$ at which
$\theta=r$, i.e.\ at which the contraction factor in the displayed chain
first dips to $r<1$.
\end{remark}

\subsection{The De~Giorgi \texorpdfstring{$L^\infty$}{L-infinity} tail
estimate}\label{ssec:Linfty}

We now prove that, for a unit-volume set, the torsion function is
quantitatively small far out, with the smallness measured by a
\emph{super-linear} power of the tail mass $\abs{\Om\setminus B_R}$. The
super-linearity (exponent $1/n>0$, or $\tfrac12-\tfrac1p>0$ when $n=2$)
is the whole point: it will let the truncation iteration converge.

\begin{lemma}[$L^\infty$ tail estimate]\label{lem:Linfty}
Let $n\ge3$. There is a dimensional constant $C_\flat=C_\flat(n)>0$ such
that for every open $\Om$ with $\abs{\Om}=\omega_n$ and every $R\ge1$,
\begin{equation}\label{eq:Linfty}
   \norm{u_\Om}_{L^\infty(\Om\setminus B_{R+1})}
   \;\le\;C_\flat\,\abs{\Om\setminus B_R}^{1/n}.
\end{equation}
For $n=2$ the same holds with the exponent $1/n$ replaced by
$\tfrac12-\tfrac1p$ for any fixed $p\in(2,\infty)$
(Remark~\textup{\ref{rem:n2}}); this weaker but still super-linear power is
all that the truncation uses.
\end{lemma}

\begin{proof}
If $\abs{\Om\setminus B_R}=0$ then $u_\Om=0$ a.e.\ on
$\Om\setminus B_R\supset\Om\setminus B_{R+1}$ and \eqref{eq:Linfty} is
trivial, so assume $\abs{\Om\setminus B_R}>0$.
%% JUSTIFICATION.  If |\Om\setminus B_R|=0 then \Om\subset B_R up to a
%% null set; since u_\Om\in H^1_0(\Om) vanishes outside \Om and \Om has
%% no mass beyond B_R, u_\Om=0 a.e.\ on \Om\setminus B_R, hence on the
%% smaller set \Om\setminus B_{R+1} (as R+1>R).  The right side of
%% \eqref{eq:Linfty} is 0 and the left side is 0, so it holds.  Below we
%% assume the tail has positive measure, so all the relative measures
%% \mu_k and normalisations a_k=\mu_k/|\Om\setminus B_R| make sense.
We carry out a De~Giorgi iteration on the annular region
$\Om\setminus B_R$. The plan: build cut-offs that switch on across a
shrinking nest of annuli, and levels rising to a target height; show
that the volume of the super-level set above the rising level, inside the
shrinking annulus, obeys a recursion of the form
\eqref{eq:DG-recursion2}; then choose the target height $M$ so large that
the seed is below the threshold \eqref{eq:DG-smallness2}, forcing the
volume to $0$ and pinning $u_\Om$ below the height.

\smallskip
\emph{Set-up: shrinking radii, cut-offs, and rising levels.}
For $k\in\mathbb N_0$ put
\[
   R_k:=R+1-2^{-k},
\]
so that $R_0=R$, the radii increase to $R_\infty=R+1$, and
$R_k\ge R\ge1$.
%% JUSTIFICATION.  R_0=R+1-1=R; as k\to\infty, 2^{-k}\to0 so R_k\to R+1;
%% k\mapsto 2^{-k} is decreasing so k\mapsto R_k is increasing, hence
%% R_k\ge R_0=R\ge1.  Thus the radii sweep the annulus B_{R+1}\setminus
%% B_R outward.
Let $\varphi_k\in C^{0,1}(\R^n)$ be the radial cut-off
\[
   \varphi_k(x):=\phi_k(\abs{x}),\qquad
   \phi_k=0\ \text{on}\ [0,R_{k-1}],\quad
   \phi_k=1\ \text{on}\ [R_k,\infty),\quad
   \phi_k\ \text{affine on}\ [R_{k-1},R_k],
\]
for $k\ge1$, and $\varphi_0\equiv1$ on $\{\abs{x}\ge R_0\}$, extended by
$0$ inside (we only use $\varphi_k$ for $k\ge1$ in the recursion). Then
$0\le\varphi_k\le1$, $\varphi_k\equiv0$ on $B_{R_{k-1}}$, $\varphi_k\equiv1$
on $\R^n\setminus B_{R_k}$, and
\[
   \abs{\nabla\varphi_k}\le\frac1{R_k-R_{k-1}}
   =\frac1{2^{-(k-1)}-2^{-k}}=2^{\,k}\le 2^{\,k+1}.
\]
%% JUSTIFICATION (the gradient bound).  \varphi_k is affine in |x| with
%% slope 1/(R_k-R_{k-1}) on the annulus B_{R_k}\setminus B_{R_{k-1}} and
%% constant elsewhere, so |\nabla\varphi_k| is that slope on the annulus
%% and 0 off it.  The width is
%%   R_k-R_{k-1}=(R+1-2^{-k})-(R+1-2^{-(k-1)})=2^{-(k-1)}-2^{-k}
%%             =2^{-k}(2-1)=2^{-k},
%% so the slope is 2^{k}; we record the slightly larger 2^{k+1} only to
%% match the algebra later (any upper bound is admissible).  The bounds
%% 0\le\varphi_k\le1 and the support facts are immediate from the
%% piecewise definition.
Choose rising levels
\[
   s_k:=M\,\abs{\Om\setminus B_R}^{1/n}\,(1-2^{-k}),\qquad k\ge0,
\]
where $M=M(n)\ge1$ is a free constant fixed at the end. Thus $s_0=0$,
$s_k\uparrow s_\infty:=M\abs{\Om\setminus B_R}^{1/n}$, and
\begin{equation}\label{eq:s-diff}
   s_{k+1}-s_k=M\abs{\Om\setminus B_R}^{1/n}\,2^{-(k+1)} .
\end{equation}
%% JUSTIFICATION.  s_0=M|\cdot|^{1/n}(1-1)=0; k\mapsto(1-2^{-k}) is
%% increasing to 1, so s_k\uparrow M|\cdot|^{1/n}=s_\infty.  The gap is
%%   s_{k+1}-s_k=M|\cdot|^{1/n}[(1-2^{-(k+1)})-(1-2^{-k})]
%%             =M|\cdot|^{1/n}(2^{-k}-2^{-(k+1)})=M|\cdot|^{1/n}2^{-(k+1)},
%% which is \eqref{eq:s-diff}.  The geometry: radii shrink the domain to
%% B_{R+1}^{\,c} and levels rise to the height s_\infty we want to prove
%% is an upper bound for u_\Om; the products (gap)^2 will furnish the
%% Chebyshev denominators.

\smallskip
\emph{Caccioppoli inequality from the truncated test function.}
Fix $k\ge1$. The function $(u_\Om-s_k)_+\in H^1_0(\Om)$ (it vanishes
where $u_\Om\le s_k$, in particular near $\partial\Om$ since $u_\Om\ge0$
and $s_k\ge0$), and $\varphi_k$ is bounded Lipschitz, so
\[
   v_k:=\varphi_k^{\,2}\,(u_\Om-s_k)_+\ \in\ H^1_0(\Om)
\]
is an admissible test function.
%% JUSTIFICATION.  (u_\Om-s_k)_+ \in H^1_0(\Om): truncation
%% w\mapsto(w-s)_+ maps H^1_0 to H^1_0 (it is 1-Lipschitz and fixes 0,
%% so preserves H^1 and the zero boundary trace), and u_\Om\in H^1_0(\Om).
%% Multiplying by the bounded Lipschitz \varphi_k^2 keeps it in H^1_0(\Om)
%% (product rule for an H^1 function times a Lipschitz function); the
%% support is inside \Om because (u_\Om-s_k)_+ already is.
Insert it into the weak formulation
$\int_\Om\nabla u_\Om\cdot\nabla v\dx=\int_\Om v\dx$. Using
$\nabla v_k=\varphi_k^2\nabla(u_\Om-s_k)_+
+2\varphi_k(u_\Om-s_k)_+\nabla\varphi_k$ and
$\nabla u_\Om\cdot\nabla(u_\Om-s_k)_+=\abs{\nabla(u_\Om-s_k)_+}^2$ (the
gradient of $u_\Om$ equals that of $(u_\Om-s_k)_+$ on $\{u_\Om>s_k\}$ and
both factors vanish elsewhere), we get
\[
   \int\varphi_k^2\abs{\nabla(u_\Om-s_k)_+}^2\dx
   +2\int\varphi_k(u_\Om-s_k)_+\,\nabla\varphi_k\cdot\nabla(u_\Om-s_k)_+\dx
   =\int\varphi_k^2(u_\Om-s_k)_+\dx .
\]
%% JUSTIFICATION.  Product rule for \nabla v_k:
%%   \nabla(\varphi_k^2 w)=\varphi_k^2\nabla w + w\nabla(\varphi_k^2)
%%     =\varphi_k^2\nabla w + 2\varphi_k w\nabla\varphi_k,  w=(u_\Om-s_k)_+.
%% Dot with \nabla u_\Om and integrate.  On \{u_\Om>s_k\}, \nabla
%% (u_\Om-s_k)_+=\nabla u_\Om, so \nabla u_\Om\cdot\nabla w=|\nabla w|^2;
%% on \{u_\Om\le s_k\}, w=0 and \nabla w=0 a.e., so both sides are 0.
%% Hence \nabla u_\Om\cdot\nabla w=|\nabla w|^2 everywhere a.e.  The
%% right side is \int v_k=\int\varphi_k^2 w by the weak equation.
The middle term is bounded by Young's inequality,
$2\abs{\varphi_k(u_\Om-s_k)_+\nabla\varphi_k\cdot\nabla(u_\Om-s_k)_+}
\le\tfrac12\varphi_k^2\abs{\nabla(u_\Om-s_k)_+}^2
+2\abs{\nabla\varphi_k}^2(u_\Om-s_k)_+^2$,
and absorbing the first piece on the left yields the Caccioppoli
estimate
\begin{equation}\label{eq:caccioppoli}
   \int\abs{\nabla\bigl(\varphi_k(u_\Om-s_k)_+\bigr)}^2\dx
   \;\le\;C_1\!\int\abs{\nabla\varphi_k}^2(u_\Om-s_k)_+^2\dx
        \;+\;C_1\!\int\varphi_k^2(u_\Om-s_k)_+\dx,
\end{equation}
with an absolute constant $C_1$ (one checks $C_1=10$ works after expanding
$\abs{\nabla(\varphi_k w)}^2\le2\varphi_k^2\abs{\nabla w}^2
+2\abs{\nabla\varphi_k}^2 w^2$ and combining with the absorbed bound; the
precise value is immaterial). Here we wrote $w=(u_\Om-s_k)_+$.
%% JUSTIFICATION (Young with the right weight).  Young: for any \eta>0,
%%   2|XY|\le \eta X^2 + \eta^{-1}Y^2.  Take
%%   X=\varphi_k|\nabla w|, Y=w|\nabla\varphi_k|, \eta=1/2:
%%   2\varphi_k w |\nabla\varphi_k||\nabla w|
%%     \le \tfrac12\varphi_k^2|\nabla w|^2 + 2 w^2|\nabla\varphi_k|^2,
%% which is the displayed bound (the dot product is \le the product of
%% norms, Cauchy--Schwarz).  Move \tfrac12\varphi_k^2|\nabla w|^2 to the
%% left of the integrated identity; the left becomes
%%   \tfrac12\int\varphi_k^2|\nabla w|^2,
%% and the surviving right side is
%%   2\int|\nabla\varphi_k|^2 w^2 + \int\varphi_k^2 w.
%% JUSTIFICATION (constant C_1=10).  Use the elementary
%%   |\nabla(\varphi_k w)|^2=|\varphi_k\nabla w+w\nabla\varphi_k|^2
%%     \le 2\varphi_k^2|\nabla w|^2 + 2 w^2|\nabla\varphi_k|^2
%% (from |a+b|^2\le2|a|^2+2|b|^2).  Bound \varphi_k^2|\nabla w|^2 by
%% twice the absorbed Caccioppoli left side, i.e.
%%   \int\varphi_k^2|\nabla w|^2\le 4\int|\nabla\varphi_k|^2 w^2
%%     + 2\int\varphi_k^2 w
%% (double the inequality 2\cdot left = \int\varphi_k^2|\nabla w|^2\le
%%  4\int|\nabla\varphi_k|^2 w^2+2\int\varphi_k^2 w).  Then
%%   \int|\nabla(\varphi_k w)|^2
%%     \le 2(4\int|\nabla\varphi_k|^2 w^2+2\int\varphi_k^2 w)
%%        +2\int|\nabla\varphi_k|^2 w^2
%%     = 10\int|\nabla\varphi_k|^2 w^2 + 4\int\varphi_k^2 w
%%     \le 10\int|\nabla\varphi_k|^2 w^2 + 10\int\varphi_k^2 w.
%% So C_1=10 is admissible; only its dimensional independence matters,
%% it folds into C_2 below.

\smallskip
\emph{Sobolev inequality.} Set
$w_k:=\varphi_k(u_\Om-s_k)_+\in H^1_0(\R^n)$ and
\[
   A_k:=\bigl\{\varphi_k(u_\Om-s_k)_+>0\bigr\}
       =\bigl\{u_\Om>s_k\bigr\}\cap\{\varphi_k>0\}
       \subset\{u_\Om>s_k\}\cap(\R^n\setminus B_{R_{k-1}}).
\]
%% JUSTIFICATION.  The product is >0 iff both factors are >0:
%% (u_\Om-s_k)_+>0 means u_\Om>s_k, and \varphi_k>0 means |x|>R_{k-1}
%% (since \varphi_k\equiv0 on B_{R_{k-1}}); hence the set equality and
%% the inclusion A_k\subset\{u_\Om>s_k\}\cap(\R^n\setminus B_{R_{k-1}}).
%% w_k has finite-measure support since \{u_\Om>s_k\}\subset\Om has
%% finite measure; w_k\in H^1_0(\R^n) extending by 0.
For $n\ge3$ the Sobolev inequality and Hölder's inequality on the support
$A_k$ give
\begin{equation}\label{eq:sobolev}
   \int w_k^2\dx
   \le\Bigl(\int w_k^{2^*}dx\Bigr)^{2/2^*}\abs{A_k}^{1-2/2^*}
   \le S_n^2\,\abs{A_k}^{2/n}\!\int\abs{\nabla w_k}^2\dx,
\end{equation}
using $1-2/2^*=2/n$.
%% JUSTIFICATION.  Hölder on the set A_k with exponents (2^*/2) and its
%% conjugate (2^*/2)'=2^*/(2^*-2): since w_k=0 off A_k,
%%   \int w_k^2=\int_{A_k}w_k^2\cdot 1
%%     \le (\int_{A_k}w_k^{2})^{?}...
%% precisely \int_{A_k}w_k^2\le(\int_{A_k}(w_k^2)^{2^*/2})^{2/2^*}
%%   |A_k|^{1-2/2^*}=(\int w_k^{2^*})^{2/2^*}|A_k|^{1-2/2^*}.  Then
%% Sobolev (\int w_k^{2^*})^{2/2^*}=\|w_k\|_{2^*}^2\le S_n^2\|\nabla
%% w_k\|_2^2=S_n^2\int|\nabla w_k|^2.  The exponent identity:
%%   1-2/2^*=1-2(n-2)/(2n)=1-(n-2)/n=2/n.
(The $n=2$ case is handled in
Remark~\ref{rem:n2}; the resulting estimate has the same shape with $2/n$
replaced by any fixed exponent in $(0,1)$, which only changes the
dimensional constants.)

Combining \eqref{eq:sobolev} with \eqref{eq:caccioppoli},
$\abs{\nabla\varphi_k}\le2^{k+1}$ (so $\abs{\nabla\varphi_k}^2\le4^{k+1}$)
and $0\le\varphi_k\le1$, we obtain
\begin{equation}\label{eq:after-sobolev}
   \int w_k^2\dx
   \le S_n^2\,C_1\,\abs{A_k}^{2/n}
   \Bigl(4^{k+1}\!\!\int_{A_k}\!(u_\Om-s_k)_+^2\dx
        +\!\!\int_{A_k}\!(u_\Om-s_k)_+\dx\Bigr),
\end{equation}
where for the second term we used
$\int\varphi_k^2(u_\Om-s_k)_+\le\int_{A_k}(u_\Om-s_k)_+$.
%% JUSTIFICATION.  Substitute \int|\nabla w_k|^2 from
%% \eqref{eq:caccioppoli} into \eqref{eq:sobolev}.  In the first
%% Caccioppoli term replace |\nabla\varphi_k|^2 by its bound 4^{k+1};
%% the integrand (u_\Om-s_k)_+^2 is supported in A_k so the integral is
%% over A_k.  In the second term, \varphi_k^2\le1 gives
%% \int\varphi_k^2(u_\Om-s_k)_+\le\int(u_\Om-s_k)_+
%% =\int_{A_k}(u_\Om-s_k)_+ (the integrand vanishes off \{u_\Om>s_k\},
%% and on A_k^c\cap\{u_\Om>s_k\}... in fact (u_\Om-s_k)_+ is supported
%% in \{u_\Om>s_k\}\supset A_k but the \varphi_k^2 factor restricts to
%% \{\varphi_k>0\}; either way the bound holds because the integrand is
%% nonnegative and we only enlarged the constant).
By the global
sup-bound \eqref{eq:global-Linfty} at $\abs{\Om}=\omega_n$ we have
$0\le(u_\Om-s_k)_+\le u_\Om\le1/(2n)\le1$ on $\Om$, hence both
$(u_\Om-s_k)_+^2\le1$ and $(u_\Om-s_k)_+\le1$ there, so each integral over
$A_k$ is at most $\abs{A_k}$. Therefore
\begin{equation}\label{eq:after-sobolev2}
   \int w_k^2\dx
   \le S_n^2\,C_1\,\bigl(4^{k+1}+1\bigr)\,\abs{A_k}^{2/n}\,\abs{A_k}
   \le C_2\,4^{k}\,\abs{A_k}^{1+2/n},
\end{equation}
with $C_2=C_2(n):=8\,S_n^2 C_1$ (using $4^{k+1}+1\le8\cdot4^{k}$).
%% JUSTIFICATION.  (u_\Om-s_k)_+\le u_\Om because s_k\ge0; and
%% u_\Om\le1/(2n)\le1 by \eqref{eq:global-Linfty} at \abs{\Om}=\omega_n.
%% So both (u_\Om-s_k)_+ and its square are \le1, and
%%   \int_{A_k}(u_\Om-s_k)_+^{(\cdot)}\le\int_{A_k}1=|A_k|.
%% Plugging into \eqref{eq:after-sobolev}:
%%   \int w_k^2\le S_n^2 C_1|A_k|^{2/n}(4^{k+1}|A_k|+|A_k|)
%%     =S_n^2 C_1(4^{k+1}+1)|A_k|^{1+2/n}.
%% Finally 4^{k+1}+1=4\cdot4^k+1\le4\cdot4^k+4^k\cdot... actually
%% 1\le4^k so 4^{k+1}+1\le4\cdot4^k+4\cdot4^k=8\cdot4^k (k\ge1 gives
%% 1\le4^k\le4\cdot4^k); set C_2=8 S_n^2 C_1.  This is the form
%% \int w_k^2\le C_2 4^k|A_k|^{1+2/n}, super-linear in |A_k|.

\smallskip
\emph{From energy to a measure recursion.} Restrict the left-hand side of
\eqref{eq:after-sobolev2} to the smaller set where the \emph{next} level
is exceeded. On $\{u_\Om>s_{k+1}\}\cap(\R^n\setminus B_{R_{k}})$ we have
$\varphi_k=1$ (since this set lies outside $B_{R_k}$) and
$u_\Om-s_k>s_{k+1}-s_k$, hence
$w_k^2=(u_\Om-s_k)_+^2\ge(s_{k+1}-s_k)^2$ there. Therefore
\begin{equation}\label{eq:chebyshev}
   (s_{k+1}-s_k)^2\,
   \bigl|\{u_\Om>s_{k+1}\}\cap(\R^n\setminus B_{R_{k}})\bigr|
   \;\le\;\int_{\R^n}w_k^2\dx .
\end{equation}
%% JUSTIFICATION (this is Chebyshev/Markov).  On the indicated set
%% \varphi_k=1 (it lies outside B_{R_k}, where \varphi_k\equiv1), so
%% w_k=(u_\Om-s_k)_+.  There u_\Om>s_{k+1}, so u_\Om-s_k>s_{k+1}-s_k>0,
%% hence w_k^2\ge(s_{k+1}-s_k)^2.  Integrating this pointwise lower
%% bound over the set and dropping the rest of \R^n (where w_k^2\ge0):
%%   (s_{k+1}-s_k)^2\,|set|\le\int_{set}w_k^2\le\int_{\R^n}w_k^2.
Introduce the (relative) super-level measures
\[
   \mu_k:=\bigl|\{u_\Om>s_k\}\cap(\R^n\setminus B_{R_{k-1}})\bigr|,
   \qquad k\ge1,
\]
which decrease in $k$ (both $s_k$ and $R_{k-1}$ increase). Since
$A_k\subset\{u_\Om>s_k\}\cap(\R^n\setminus B_{R_{k-1}})$ we have
$\abs{A_k}\le\mu_k$, and the left-hand set in \eqref{eq:chebyshev} is
exactly $\{u_\Om>s_{k+1}\}\cap(\R^n\setminus B_{R_{(k+1)-1}})$, of
measure $\mu_{k+1}$.
%% JUSTIFICATION.  As k grows, s_k\uparrow and R_{k-1}\uparrow, so both
%% defining conditions \{u_\Om>s_k\} and (\R^n\setminus B_{R_{k-1}})
%% shrink; hence \mu_k is nonincreasing.  |A_k|\le\mu_k by the inclusion
%% recorded above.  And R_{(k+1)-1}=R_k, so the set in \eqref{eq:chebyshev}
%% is precisely \{u_\Om>s_{k+1}\}\cap(\R^n\setminus B_{R_k})=\mu_{k+1}'s
%% set.
Combining \eqref{eq:chebyshev},
\eqref{eq:after-sobolev2} and $\abs{A_k}\le\mu_k$,
\[
   (s_{k+1}-s_k)^2\,\mu_{k+1}
   \;\le\;C_2\,4^{k}\,\mu_k^{1+2/n}.
\]
%% JUSTIFICATION.  Chain: (s_{k+1}-s_k)^2\mu_{k+1}\le\int w_k^2
%% \le C_2 4^k|A_k|^{1+2/n}\le C_2 4^k\mu_k^{1+2/n}, the last step since
%% |A_k|\le\mu_k and t\mapsto t^{1+2/n} is increasing on [0,\infty).
Insert the gap \eqref{eq:s-diff},
$(s_{k+1}-s_k)^2=M^2\abs{\Om\setminus B_R}^{2/n}4^{-(k+1)}$:
\begin{equation}\label{eq:mu-recursion}
   \mu_{k+1}\;\le\;
   \frac{C_2\,4^{k}}{M^2\abs{\Om\setminus B_R}^{2/n}4^{-(k+1)}}\,
   \mu_k^{1+2/n}
   \;=\;\frac{4\,C_2}{M^2\abs{\Om\setminus B_R}^{2/n}}\,16^{\,k}\,
   \mu_k^{1+2/n}.
\end{equation}
%% JUSTIFICATION.  From \eqref{eq:s-diff},
%% (s_{k+1}-s_k)^2=M^2|\cdot|^{2/n}(2^{-(k+1)})^2=M^2|\cdot|^{2/n}4^{-(k+1)}.
%% Divide the previous display by it; the k-power becomes
%% 4^k/4^{-(k+1)}=4^{k+(k+1)}=4^{2k+1}=4\cdot16^k.  So the prefactor is
%% (C_2\cdot4)/(M^2|\cdot|^{2/n}) times 16^k, as displayed.

\smallskip
\emph{Normalisation and the iteration lemma.} Set
$a_k:=\mu_k/\abs{\Om\setminus B_R}$. Since
$\mu_k\le\abs{\R^n\setminus B_{R_{k-1}}\cap\{u_\Om>s_k\}}
\le\abs{\Om\setminus B_R}$ (because $R_{k-1}\ge R_0=R$ and
$\{u_\Om>s_k\}\subset\Om$), we have $0\le a_k\le1$; in particular
$a_1\le1$.
%% JUSTIFICATION.  \mu_k's set is \{u_\Om>s_k\}\cap(\R^n\setminus
%% B_{R_{k-1}})\subset \Om\cap(\R^n\setminus B_R)=\Om\setminus B_R, using
%% \{u_\Om>s_k\}\subset\{u_\Om>0\}\subset\Om and B_{R_{k-1}}\supset B_R
%% (as R_{k-1}\ge R).  Hence \mu_k\le|\Om\setminus B_R| and a_k\in[0,1].
Dividing \eqref{eq:mu-recursion} by
$\abs{\Om\setminus B_R}$ and using
$\mu_k^{1+2/n}=\abs{\Om\setminus B_R}^{1+2/n}a_k^{1+2/n}$,
\[
   a_{k+1}
   \le\frac{4C_2}{M^2\abs{\Om\setminus B_R}^{2/n}}\,16^k\,
   \frac{\abs{\Om\setminus B_R}^{1+2/n}}{\abs{\Om\setminus B_R}}
   \,a_k^{1+2/n}
   =\frac{4C_2}{M^2}\,16^{k}\,a_k^{1+2/n},
\]
valid for all $k\ge1$.
%% JUSTIFICATION (the powers of |\Om\setminus B_R| cancel).  Write
%% T:=|\Om\setminus B_R|.  Then \mu_{k+1}=T a_{k+1},
%% \mu_k^{1+2/n}=T^{1+2/n}a_k^{1+2/n}.  Divide \eqref{eq:mu-recursion}
%% by T:
%%   a_{k+1}=\mu_{k+1}/T\le \frac{4C_2}{M^2 T^{2/n}}16^k
%%        \frac{T^{1+2/n}}{T}a_k^{1+2/n}.
%% The T-exponent is -2/n+(1+2/n)-1=0, so all T's cancel, leaving
%%   a_{k+1}\le(4C_2/M^2)16^k a_k^{1+2/n}.  This is the crucial
%% scale invariance: the recursion no longer sees the (unknown) tail mass.
To put this in the exact form
\eqref{eq:DG-recursion2} (indexed from $0$) we shift: set
$\widetilde a_j:=a_{j+1}$ for $j\ge0$. Then for every $j\ge0$,
\[
   \widetilde a_{j+1}=a_{j+2}
   \le\frac{4C_2}{M^2}\,16^{\,j+1}\,a_{j+1}^{1+2/n}
   =\underbrace{\frac{64\,C_2}{M^2}}_{=:C}\,16^{\,j}\,
     \widetilde a_j^{\,1+2/n},
\]
which is \eqref{eq:DG-recursion2} with
\[
   C:=\frac{64\,C_2}{M^2},\qquad b:=16,\qquad\beta:=\frac2n .
\]
%% JUSTIFICATION (the index shift).  The recursion a_{k+1}\le(4C_2/M^2)
%% 16^k a_k^{1+2/n} holds for k\ge1.  Set \widetilde a_j=a_{j+1}, j\ge0,
%% so the seed is \widetilde a_0=a_1\le1 (which we control), not a_0.
%% Apply the recursion at k=j+1\ge1:
%%   \widetilde a_{j+1}=a_{j+2}\le(4C_2/M^2)16^{j+1}a_{j+1}^{1+2/n}
%%     =(4C_2/M^2)\cdot16\cdot16^j\widetilde a_j^{1+2/n}
%%     =(64C_2/M^2)16^j\widetilde a_j^{1+2/n}.
%% This matches \eqref{eq:DG-recursion2} with C=64C_2/M^2, b=16,
%% \beta=2/n.  (We absorbed the extra factor 16 from 16^{j+1} into C.)
Choose $M=M(n)$ so large that the smallness condition
\eqref{eq:DG-smallness2} holds at the starting value
$\widetilde a_0=a_1\le1$. Since $a_1\le1$ it suffices to require
\[
   1\;\le\;C^{-1/\beta}\,b^{-1/\beta^2}
   =\Bigl(\tfrac{64C_2}{M^2}\Bigr)^{-n/2}16^{-n^2/4},
   \quad\text{i.e.}\quad
   M^2\;\ge\;64\,C_2\,16^{\,n^2/4\cdot(2/n)}=64\,C_2\,16^{\,n/2},
\]
which holds for $M:=8\,(C_2)^{1/2}4^{\,n/2}$ (a dimensional constant,
since $C_2=C_2(n)$). With this $M$, Lemma~\ref{lem:DG-iteration2} gives
$\widetilde a_j\to0$, i.e.\ $a_k\to0$ and hence $\mu_k\to0$.
%% JUSTIFICATION (the algebra of the threshold).  With \beta=2/n,
%% 1/\beta=n/2 and 1/\beta^2=n^2/4.  The threshold \eqref{eq:DG-smallness2}
%% is \widetilde a_0\le C^{-1/\beta}b^{-1/\beta^2}; since \widetilde a_0
%% =a_1\le1, it suffices that 1\le C^{-n/2}16^{-n^2/4}, i.e.
%%   C^{n/2}16^{n^2/4}\le1?  No --- rearrange properly:
%%   1\le C^{-n/2}16^{-n^2/4}  \iff  C^{n/2}\cdot16^{n^2/4}\le1?
%% That cannot hold for large constants; instead read it as a LOWER
%% bound on M.  Raise to the power 2/n:
%%   1\le C^{-1}16^{-(n^2/4)(2/n)}=C^{-1}16^{-n/2},  i.e.  C\le16^{-n/2}?
%% Again rearrange in M: C=64C_2/M^2, so C\,16^{n/2}\le1 means
%%   64C_2 16^{n/2}/M^2\le1, i.e.  M^2\ge64C_2 16^{n/2}.
%% Thus the requirement is M^2\ge64C_2 16^{n/2}=64C_2 4^{n}, satisfied by
%%   M=8\sqrt{C_2}\,4^{n/2}  (then M^2=64 C_2 4^n=64C_2 16^{n/2}).
%% [The mid-line scratch above only shows the bookkeeping of exponents;
%% the binding statement is M^2\ge64C_2 16^{n/2}.]  M depends only on n
%% since C_2=C_2(n).  Lemma~\ref{lem:DG-iteration2} now applies to
%% (\widetilde a_j) with this C,b,\beta and seed \le threshold, giving
%% \widetilde a_j\to0.  Since a_k=\widetilde a_{k-1} for k\ge1 and \mu_k
%% =T a_k, also \mu_k\to0.

\smallskip
\emph{Conclusion.} Letting $k\to\infty$ in $\mu_k\to0$ and using
$s_k\uparrow s_\infty=M\abs{\Om\setminus B_R}^{1/n}$ and
$R_{k-1}\uparrow R+1$, monotone convergence yields
\[
   \bigl|\{u_\Om>s_\infty\}\cap(\R^n\setminus B_{R+1})\bigr|
   =\lim_{k\to\infty}\mu_k=0 .
\]
%% JUSTIFICATION.  The sets E_k:=\{u_\Om>s_k\}\cap(\R^n\setminus
%% B_{R_{k-1}}) decrease in k, and their intersection contains (indeed
%% equals, up to the increasing-union subtlety) the limiting set
%% \{u_\Om>s_\infty\}\cap(\R^n\setminus B_{R+1}): if u_\Om(x)>s_\infty and
%% |x|>R+1 then for every k, u_\Om(x)>s_\infty\ge s_k and |x|>R+1>R_{k-1},
%% so x\in E_k.  Hence
%%   |\{u_\Om>s_\infty\}\cap(\R^n\setminus B_{R+1})|\le\inf_k\mu_k
%%    =\lim_k\mu_k=0
%% by monotone convergence (continuity of measure along a decreasing
%% sequence of finite-measure sets).
Hence $u_\Om\le s_\infty=M\abs{\Om\setminus B_R}^{1/n}$ a.e.\ on
$\Om\setminus B_{R+1}$. Since $u_\Om\ge0$, this is
\eqref{eq:Linfty} with $C_\flat:=M=8\,(C_2)^{1/2}4^{n/2}$.
%% JUSTIFICATION.  The null set \{u_\Om>s_\infty\}\cap(\R^n\setminus
%% B_{R+1}) being empty up to measure zero means u_\Om\le s_\infty a.e.
%% on \R^n\setminus B_{R+1}, hence a.e. on \Om\setminus B_{R+1}.  As
%% u_\Om\ge0 there, \|u_\Om\|_{L^\infty(\Om\setminus B_{R+1})}\le
%% s_\infty=M|\Om\setminus B_R|^{1/n}, which is \eqref{eq:Linfty} with
%% C_\flat=M (dimensional).
\end{proof}

\begin{remark}[The case $n=2$]\label{rem:n2}
For $n=2$ the embedding $H^1_0\hookrightarrow L^{2^*}$ degenerates
(formally $2^*=2n/(n-2)=\infty$), but the
Gagliardo--Nirenberg--Sobolev inequality gives, for any fixed
$p\in(2,\infty)$,
$\norm{w}_{L^{p}(\R^2)}\le S_{2,p}\norm{\nabla w}_{L^2(\R^2)}$ for
$w\in H^1_0(\R^2)$ with finite-measure support. Repeating the Hölder step
with exponent $p$ replaces $2/n=1$ by $1-2/p\in(0,1)$, so
\eqref{eq:mu-recursion} becomes
$\mu_{k+1}\le C\,16^k\,\mu_k^{1+\beta}$ with $\beta:=1-2/p>0$, and
Lemma~\ref{lem:DG-iteration2} applies. Tracking the level scale: with
$s_\infty=M\abs{\Om\setminus B_R}^{\alpha}$, self-consistency of the
recursion forces $\alpha=\tfrac{p-2}{2p}=\tfrac12-\tfrac1p$ (for $n\ge3$
the analogous bookkeeping gives $\alpha=1/n$, and the two agree only there).
Hence, for $n=2$, the $L^\infty$ tail estimate holds in the slightly weaker
form
\[
   \norm{u_\Om}_{L^\infty(\Om\setminus B_{R+1})}
   \;\le\;C(p)\,\abs{\Om\setminus B_R}^{\,\frac12-\frac1p},
   \qquad\text{any fixed }p\in(2,\infty),
\]
in place of the exponent $1/n$ in \eqref{eq:Linfty}. Fixing once and for
all $p=4$, so that $\tfrac12-\tfrac1p=\tfrac14>0$, this is a genuine
super-linear power of the tail mass, which is all the truncation of
\S\ref{ssec:truncation} uses: the energy-comparison
\eqref{eq:energy-comp} and the tail recursion \eqref{eq:bk2-bound} hold
with the exponent $1+\tfrac1n$ replaced by $1+\kappa$, $\kappa=\tfrac14$,
and the geometric iteration \eqref{eq:bk2-self}--\eqref{eq:bodd-decay} goes
through for any $\kappa>0$ with $n$-dependent constants. Consequently the
global Theorem~\ref{thm:SV-global} holds for all $n\ge2$; only the
auxiliary exponent in \eqref{eq:Linfty} (and the resulting dimensional
constants) changes at $n=2$.
%% JUSTIFICATION (the self-consistency of \alpha).  In the n\ge3 proof the
%% level scale was s_\infty=M|\cdot|^{1/n}, and this exponent 1/n was
%% forced by exactly the cancellation of powers of T=|\Om\setminus B_R|
%% in the normalisation step: there the Sobolev/Hölder exponent is 2/n
%% and the gap-squared contributes |\cdot|^{2\alpha} in the denominator;
%% the T-exponent vanished precisely because 2\alpha=2/n and the recursion
%% closed.  For p the Hölder exponent becomes 1-2/p=\beta, the Sobolev
%% step produces |A_k|^{\beta} (not |A_k|^{2/n}), and the same T-cancellation
%% now demands 2\alpha=\beta=1-2/p, i.e. \alpha=(1-2/p)/2=1/2-1/p.  At p=4
%% this is 1/4.  Everything downstream uses only that \alpha>0 and that
%% the recursion is super-linear, both of which hold; the constants pick
%% up a dependence on p (hence on the fixed choice p=4), still dimensional.
\end{remark}

\subsection{A suboptimal stability seed}\label{ssec:seed}

The truncation iteration needs one quantitative input that the De~Giorgi
machinery does not by itself supply: a crude (far from sharp) control of the
asymmetry by the deficit, used only to make the tail mass small at the first
radius and to bound the number of truncation steps. We prove it here from the
level-set identity already established, so that the present section introduces
no external stability input. The mechanism is: the deficit identity plus the
quantitative isoperimetric inequality control a weighted integral of the
asymmetries of the superlevel sets; the profile-gap identity shows that this
weight has positive mass on the outer collar; and the superlevel transfer
shows that on that collar the superlevel sets inherit the asymmetry of
$\Om$ itself. The fourth power (in fact a third power) is the price of the
crude collar transfer; the sharp second power is the main theorem.

\begin{lemma}[Suboptimal stability seed]\label{lem:seed}
There is a dimensional constant $C_\dagger=C_\dagger(n)>0$ such that every
open $\Om\subset\R^n$ of finite measure satisfies
\begin{equation}\label{eq:seed}
   D(\Om)\;\ge\;\frac{1}{C_\dagger}\,\Fr(\Om)^{4}.
\end{equation}
Equivalently, $\Fr(\Om)\le\bigl(C_\dagger D(\Om)\bigr)^{1/4}$.
\end{lemma}

\begin{proof}
$D$ and $\Fr$ are scale invariant, so we may assume $\abs\Om=\omega_n$; then
$D(\Om)=\omega_n^{-(n+2)/n}\dT(\Om)$, $B^*=B_1$. Write $A:=\Fr(\Om)\in[0,2)$,
and recall $E_\rho=\{u>t(\rho)\}$, $\abs{E_\rho}=\omega_n\rho^n$,
$a(\rho)=-t'(\rho)$. Three facts proved earlier are used; each holds for any
open set of finite measure, since its proof does not use $\Om\subset B_R$:
\begin{itemize}
\item[(a)] the exact deficit identity (Theorem~\ref{thm:id-main}),
$2\dT(\Om)=\int_0^{\norm u_\infty}(D_H+D_I)/(n^2\omega_n^{2/n}m^{1-2/n})\,dt$,
with $D_H,D_I\ge0$;
\item[(b)] the profile-gap identity \eqref{eq:profile-gap-rho}:
$0\le a(\rho)\le\rho/n$ and
$\int_0^1\rho^n\bigl(\tfrac\rho n-a(\rho)\bigr)\,d\rho=\tfrac2{\omega_n}\dT(\Om)$;
\item[(c)] the superlevel transfer (Lemma~\ref{lem:superlevel-transfer}):
$A\le\Fr(E_\rho)+2(1-\rho^n)$ for every $\rho\in(0,1]$.
\end{itemize}

\emph{Step 1 (a level-set lower bound for $D_I$).}
Fix a finite-perimeter level $\rho$. Fusco--Julin
(Theorem~\ref{thm:FJ-strong}) applied to $E_\rho$ with $r_E=\rho$, after
discarding the nonnegative boundary-oscillation summand and the outer square
in \eqref{eq:FJ-exact}, reads
$\abs{E_\rho\Delta B_\rho(z)}^2/\rho^{2n}\le C_{\rm FJ}(n)
(\Per(E_\rho)-\Per(B_\rho))/\rho^{n-1}$ for the optimal centre $z$, i.e.
$\abs{E_\rho\Delta B_\rho(z)}^2\le C_{\rm FJ}(n)\rho^{n+1}
(\Per(E_\rho)-\Per(B_\rho))$.
%% JUSTIFICATION. Both summands in \eqref{eq:FJ-exact} are \ge0, so each is
%% \le the right side; drop the second and the square, then multiply by
%% r_E^{2n}=\rho^{2n} and divide r_E^{n-1}=\rho^{n-1}.
Now $\Fr(E_\rho)\le\abs{E_\rho\Delta B_\rho(z)}/\abs{E_\rho}$ (the Fraenkel
asymmetry is an infimum over balls of volume $\abs{E_\rho}$, and $B_\rho(z)$
is one such), and
$\Per(E_\rho)-\Per(B_\rho)=D_I(t(\rho))/(\Per(E_\rho)+\Per(B_\rho))\le
D_I(t(\rho))/\Per(B_\rho)$. With $\abs{E_\rho}=\omega_n\rho^n$ and
$\Per(B_\rho)=n\omega_n\rho^{n-1}$,
\[
   \Fr(E_\rho)^2\,\omega_n^2\rho^{2n}\le\abs{E_\rho\Delta B_\rho(z)}^2
   \le C_{\rm FJ}(n)\rho^{n+1}\frac{D_I(t(\rho))}{n\omega_n\rho^{n-1}}
   =\frac{C_{\rm FJ}(n)}{n\omega_n}\rho^2 D_I(t(\rho)),
\]
which rearranges to
\begin{equation}\label{eq:seed-FJ}
   D_I(t(\rho))\;\ge\;c_1(n)\,\rho^{\,2n-2}\,\Fr(E_\rho)^2,
   \qquad c_1(n):=\frac{n\,\omega_n^{3}}{C_{\rm FJ}(n)} .
\end{equation}

\emph{Step 2 (restrict to the outer collar $[\rho_A,1]$).}
Put $\rho_A:=(1-A/4)^{1/n}$. For $\rho\in[\rho_A,1]$,
$1-\rho^n\le1-\rho_A^n=A/4$, so (c) gives
$\Fr(E_\rho)\ge A-2(1-\rho^n)\ge A/2$, and $\rho^n\ge\rho_A^n=1-A/4>1/2$.
In (a) drop $D_H\ge0$ and change variables $t=t(\rho)$
($dt=a(\rho)\,d\rho$, $m=\omega_n\rho^n$); keeping only the nonnegative
integrand over $[\rho_A,1]$ and inserting \eqref{eq:seed-FJ},
\[
   2\dT(\Om)\ \ge\
   \int_{\rho_A}^1\frac{D_I(t(\rho))}{n^2\omega_n^{2/n}(\omega_n\rho^n)^{1-2/n}}\,a(\rho)\,d\rho
   \ \ge\ c_2(n)\int_{\rho_A}^1\rho^n\,\Fr(E_\rho)^2\,a(\rho)\,d\rho,
\]
where $c_2(n)=c_1(n)/(n^2\omega_n)$.
%% JUSTIFICATION (constant). (\omega_n\rho^n)^{1-2/n}=\omega_n^{1-2/n}\rho^{n-2};
%% so n^2\omega_n^{2/n}(\omega_n\rho^n)^{1-2/n}=n^2\omega_n\rho^{n-2}, and
%% \rho^{2n-2}/\rho^{n-2}=\rho^n. Hence c_1\rho^{2n-2}/(n^2\omega_n\rho^{n-2})
%% =(c_1/(n^2\omega_n))\rho^n.
With $\Fr(E_\rho)\ge A/2$ on $[\rho_A,1]$,
\begin{equation}\label{eq:seed-collar}
   2\dT(\Om)\ \ge\ \frac{c_2(n)\,A^2}{4}\int_{\rho_A}^1\rho^n a(\rho)\,d\rho .
\end{equation}

\emph{Step 3 (the profile gap bounds the collar mass below).}
By (b) and $\tfrac\rho n-a\ge0$,
\[
   \int_{\rho_A}^1\rho^n a\,d\rho
   =\int_{\rho_A}^1\frac{\rho^{n+1}}{n}\,d\rho
     -\int_{\rho_A}^1\rho^n\Bigl(\frac\rho n-a\Bigr)d\rho
   \ \ge\ \frac{1-\rho_A^{\,n+2}}{n(n+2)}-\frac2{\omega_n}\dT(\Om),
\]
the inequality because the subtracted integral is at most the full integral
$\int_0^1\rho^n(\tfrac\rho n-a)=\tfrac2{\omega_n}\dT(\Om)$. Since $\rho_A\le1$,
$\rho_A^{\,n+2}\le\rho_A^n=1-A/4$, so $1-\rho_A^{\,n+2}\ge A/4$ and
\begin{equation}\label{eq:seed-mass}
   \int_{\rho_A}^1\rho^n a\,d\rho\ \ge\ \frac{A}{4n(n+2)}-\frac2{\omega_n}\dT(\Om).
\end{equation}

\emph{Step 4 (conclusion).} Insert \eqref{eq:seed-mass} into
\eqref{eq:seed-collar}:
\[
   2\dT(\Om)\ \ge\ \frac{c_2(n)A^2}{4}
   \Bigl(\frac{A}{4n(n+2)}-\frac2{\omega_n}\dT(\Om)\Bigr)
   = c_3(n)A^3-c_4(n)A^2\dT(\Om),
\]
$c_3(n)=c_2(n)/(16n(n+2))$, $c_4(n)=c_2(n)/(2\omega_n)$. Therefore
$\dT(\Om)(2+c_4(n)A^2)\ge c_3(n)A^3$, and since $A<2$,
$\dT(\Om)\ge c_3(n)A^3/(2+4c_4(n))=:c_5(n)A^3$. Finally $A<2$ gives
$A^4\le2A^3$, hence $\dT(\Om)\ge\tfrac12 c_5(n)A^4$; multiplying by
$\omega_n^{-(n+2)/n}$ yields \eqref{eq:seed} with
$C_\dagger(n)=2\omega_n^{(n+2)/n}/c_5(n)$.
\end{proof}

\begin{remark}[The seed is internal]\label{rem:seed-status}
The only external ingredient is the Fusco--Julin quantitative isoperimetric
inequality (Theorem~\ref{thm:FJ-strong}), already used in the manuscript;
the deficit identity, the profile-gap identity, and the superlevel transfer
are all proved here. So the bounded reduction rests on no stability estimate
beyond those the paper establishes. The exponent obtained ($3$, hence the
stated $4$) is far from the sharp exponent $2$ of the main theorem, but only
a fixed positive power of $D$ is used below. A suboptimal torsional stability
bound of this kind is classical; cf.\ Fusco--Maggi--Pratelli \cite{FMP}.
\end{remark}

\subsection{The truncation construction}\label{ssec:truncation}

We now build the bounded surrogate. The estimate driving the construction
is an \emph{energy-comparison inequality}: cutting $\Om$ off at radius
$\sim k$ costs only a higher-order amount of energy, controlled by the
tail mass through Lemma~\ref{lem:Linfty}. The logic of the proof is a
feedback loop: the energy comparison (Step~1) becomes, via Saint--Venant,
a recursion bounding the tail mass two radii out by the tail mass here
plus the deficit (Step~2); a stopping index $K$ marks where the tail
first drops below the deficit scale (Step~3); the recursion is shown to
force $K$ to be dimensionally bounded (the heart of the matter); finally
we cut at $K+3$, rescale to unit volume, and verify the four advertised
properties (Steps~5--7).

\begin{lemma}[Bounded truncation]\label{lem:truncation}
There exist dimensional constants $C_\sharp=C_\sharp(n)>0$,
$\delta=\delta(n)\in(0,1]$ and $d=d(n)\ge1$ such that for every open
$\Om$ with $\abs{\Om}=\omega_n$ and $D(\Om)\le\delta$ there is an open set
$\Otilde$ with
\begin{align}
   \abs{\Otilde}\;&=\;\omega_n,\label{eq:trunc-vol}\\
   \diam(\Otilde)\;&\le\;d,\label{eq:trunc-diam}\\
   D(\Otilde)\;&\le\;C_\sharp\,D(\Om),\label{eq:trunc-D}\\
   \Fr(\Om)\;&\le\;\Fr(\Otilde)+C_\sharp\,D(\Om).\label{eq:trunc-A}
\end{align}
\end{lemma}

\begin{proof}
Normalise so that the ball realising the asymmetry of $\Om$ is $B_1$
(translate $\Om$; both $\Fr$ and $D$ are translation invariant), i.e.
$\Fr(\Om)=\abs{\Om\,\triangle\,B_1}/\omega_n$.
%% JUSTIFICATION.  \Fr is an infimum over translates of B_1, and the
%% infimum is attained for a fixed-volume open set (a standard
%% compactness of the centre x_0, or one works with a near-optimal ball;
%% either suffices).  Translating \Om so the optimal ball is centred at 0
%% changes neither \Fr nor D (both translation invariant: \Fr by its
%% definition, D because E and \abs{\cdot} are translation invariant).
We will repeatedly use the tail-mass notation
\begin{equation}\label{eq:bk-def}
   b_k:=\frac{\abs{\Om\setminus B_k}}{\omega_n},\qquad k\ge1,
\end{equation}
which is non-increasing in $k$ and satisfies $0\le b_k\le1$.
%% JUSTIFICATION.  k\mapsto B_k increases, so \Om\setminus B_k decreases,
%% so |\Om\setminus B_k| and hence b_k are nonincreasing.  And
%% |\Om\setminus B_k|\le|\Om|=\omega_n gives b_k\le1; b_k\ge0 trivially.

\smallskip
\emph{Step 0: the tail mass is small.}
Since $B_1$ realises the asymmetry,
$\abs{\Om\setminus B_1}\le\abs{\Om\,\triangle\,B_1}=\omega_n\Fr(\Om)$.
%% JUSTIFICATION.  \Om\setminus B_1\subset\Om\triangle B_1 (the symmetric
%% difference contains everything in \Om but not in B_1), so the measures
%% compare; and |\Om\triangle B_1|=\omega_n\Fr(\Om) by the normalisation.
Combined with the suboptimal seed Lemma~\ref{lem:seed},
\begin{equation}\label{eq:b1-small}
   b_1\;\le\;\Fr(\Om)\;\le\;\bigl(C_\dagger\,D(\Om)\bigr)^{1/4}.
\end{equation}
%% JUSTIFICATION.  Divide the previous display by \omega_n: b_1\le\Fr(\Om);
%% then \Fr(\Om)\le(C_\dagger D(\Om))^{1/4} by Lemma~\ref{lem:seed}.
In particular, since $b_k\le b_1$ for all $k\ge1$, every tail mass is
$\le(C_\dagger D(\Om))^{1/4}$, which can be made smaller than any
prescribed dimensional threshold by shrinking $\delta(n)$ (recall
$D(\Om)\le\delta(n)$). This is the only use of Lemma~\ref{lem:seed} in
the construction.
%% JUSTIFICATION.  b_k\le b_1 by monotonicity; and (C_\dagger D)^{1/4}\le
%% (C_\dagger\delta)^{1/4}\to0 as \delta\to0, so any threshold is met by
%% taking \delta(n) small.  We invoke this twice below
%% (\eqref{eq:seed-smallness} and the L-large step).

\smallskip
\emph{Step 1: the energy-comparison inequality.}
For $k\ge1$ let $\varphi_k$ be the Lipschitz cut-off
\[
   \varphi_k(x):=\min\bigl\{1,(k+2-\abs{x})_+\bigr\},
\]
so $\varphi_k\equiv1$ on $B_{k+1}$, $\varphi_k\equiv0$ outside $B_{k+2}$,
$0\le\varphi_k\le1$ and $\abs{\nabla\varphi_k}\le1$ everywhere (it is
$1$ only on the annulus $B_{k+2}\setminus B_{k+1}$).
%% JUSTIFICATION.  For |x|\le k+1, k+2-|x|\ge1 so the min is 1.  For
%% |x|\ge k+2, (k+2-|x|)_+=0 so \varphi_k=0.  On k+1\le|x|\le k+2,
%% \varphi_k=k+2-|x| is affine in |x| with |\nabla\varphi_k|=
%% |\nabla(k+2-|x|)|=|{-\hat x}|=1.  Elsewhere \varphi_k is locally
%% constant, gradient 0.  Hence 0\le\varphi_k\le1 and |\nabla\varphi_k|
%% \le1 with equality only on the annulus.
Then
$u_k:=\varphi_k\,u_\Om\in H^1_0(\Om\cap B_{k+2})$ is admissible for the
variational problem defining $E(\Om\cap B_{k+2})$, so
$E(\Om\cap B_{k+2})\le J_{\Om\cap B_{k+2}}(u_k)
=\tfrac12\int\abs{\nabla u_k}^2-\int u_k$.
%% JUSTIFICATION.  u_k=\varphi_k u_\Om is in H^1_0 of \Om (product of
%% H^1_0 and bounded Lipschitz) and vanishes outside B_{k+2} (where
%% \varphi_k=0), hence u_k\in H^1_0(\Om\cap B_{k+2}).  E(\Om\cap B_{k+2})
%% is the MINIMUM of J over that space, so it is \le the value at the
%% admissible u_k.
We use the pointwise algebraic
identity
\[
   \abs{\nabla(\varphi_k u_\Om)}^2
   =\nabla u_\Om\cdot\nabla(u_\Om\varphi_k^2)
   +\abs{\nabla\varphi_k}^2u_\Om^2,
\]
which is verified by expanding both sides
($\nabla(u_\Om\varphi_k^2)=\varphi_k^2\nabla u_\Om
+2\varphi_k u_\Om\nabla\varphi_k$, so the right-hand side equals
$\varphi_k^2\abs{\nabla u_\Om}^2+2\varphi_k u_\Om\nabla u_\Om\cdot
\nabla\varphi_k+\abs{\nabla\varphi_k}^2 u_\Om^2
=\abs{\varphi_k\nabla u_\Om+u_\Om\nabla\varphi_k}^2$).
%% JUSTIFICATION.  Left side:
%%   |\nabla(\varphi_k u_\Om)|^2=|\varphi_k\nabla u_\Om+u_\Om\nabla
%%   \varphi_k|^2=\varphi_k^2|\nabla u_\Om|^2+2\varphi_k u_\Om\nabla
%%   u_\Om\cdot\nabla\varphi_k+u_\Om^2|\nabla\varphi_k|^2.
%% Right side: \nabla u_\Om\cdot(\varphi_k^2\nabla u_\Om+2\varphi_k u_\Om
%%   \nabla\varphi_k)+|\nabla\varphi_k|^2 u_\Om^2
%%   =\varphi_k^2|\nabla u_\Om|^2+2\varphi_k u_\Om\nabla u_\Om\cdot
%%   \nabla\varphi_k+|\nabla\varphi_k|^2 u_\Om^2.
%% Both sides coincide.  This is the standard "energy splitting" that
%% trades |\nabla(\varphi u)|^2 for a term testable by the weak equation
%% plus a controllable gradient-of-cutoff term.
Integrating and
applying the weak equation
$\int\nabla u_\Om\cdot\nabla(u_\Om\varphi_k^2)=\int u_\Om\varphi_k^2$
with the admissible test function $u_\Om\varphi_k^2\in H^1_0(\Om)$, we
get
\[
   J_{\Om\cap B_{k+2}}(u_k)
   =\tfrac12\!\int u_\Om\varphi_k^2
   +\tfrac12\!\int\abs{\nabla\varphi_k}^2u_\Om^2
   -\int\varphi_k u_\Om .
\]
%% JUSTIFICATION.  J(u_k)=\tfrac12\int|\nabla u_k|^2-\int u_k.  By the
%% identity, \tfrac12\int|\nabla u_k|^2=\tfrac12\int\nabla u_\Om\cdot
%% \nabla(u_\Om\varphi_k^2)+\tfrac12\int|\nabla\varphi_k|^2 u_\Om^2;
%% the first integral is \tfrac12\int u_\Om\varphi_k^2 by the weak
%% equation with test function u_\Om\varphi_k^2 (admissible: H^1_0 times
%% bounded Lipschitz).  And \int u_k=\int\varphi_k u_\Om.  Collect.
Now use $\tfrac12\varphi_k^2-\varphi_k=-\tfrac12(1-(1-\varphi_k)^2)
=-\tfrac12+\tfrac12(1-\varphi_k)^2$ to write
$\tfrac12\int u_\Om\varphi_k^2-\int\varphi_k u_\Om
=-\tfrac12\int u_\Om+\tfrac12\int(1-\varphi_k)^2u_\Om
=E(\Om)+\tfrac12\int(1-\varphi_k)^2u_\Om$, where we used
$E(\Om)=-\tfrac12\int u_\Om$ from \eqref{eq:energy-identities}.
%% JUSTIFICATION (the algebraic identity).  Expand
%%   1-(1-\varphi_k)^2=1-(1-2\varphi_k+\varphi_k^2)=2\varphi_k-\varphi_k^2,
%% so \tfrac12(1-(1-\varphi_k)^2)=\varphi_k-\tfrac12\varphi_k^2, hence
%%   \tfrac12\varphi_k^2-\varphi_k=-(\varphi_k-\tfrac12\varphi_k^2)
%%     =-\tfrac12(1-(1-\varphi_k)^2)=-\tfrac12+\tfrac12(1-\varphi_k)^2.
%% Multiply by u_\Om and integrate:
%%   \tfrac12\int\varphi_k^2 u_\Om-\int\varphi_k u_\Om
%%     =-\tfrac12\int u_\Om+\tfrac12\int(1-\varphi_k)^2 u_\Om
%%     =E(\Om)+\tfrac12\int(1-\varphi_k)^2 u_\Om,
%% using E(\Om)=-\tfrac12\int u_\Om.
Therefore
\begin{equation}\label{eq:energy-exact}
   E(\Om\cap B_{k+2})\le J_{\Om\cap B_{k+2}}(u_k)
   =E(\Om)+\tfrac12\int(1-\varphi_k)^2u_\Om
          +\tfrac12\int\abs{\nabla\varphi_k}^2u_\Om^2 .
\end{equation}
%% JUSTIFICATION.  Substitute the last identity into the J-expression:
%% J(u_k)=[E(\Om)+\tfrac12\int(1-\varphi_k)^2u_\Om]
%%        +\tfrac12\int|\nabla\varphi_k|^2 u_\Om^2.  This is an EXACT
%% identity for J(u_k); the inequality is the variational \le from above.
Both correction integrals are supported in $\R^n\setminus B_{k+1}$:
indeed $1-\varphi_k=0$ on $B_{k+1}$ and $\nabla\varphi_k=0$ on
$B_{k+1}\cup(\R^n\setminus B_{k+2})$. Hence, with $0\le1-\varphi_k\le1$
and $\abs{\nabla\varphi_k}\le1$,
\begin{equation}\label{eq:energy-comp-pre}
   E(\Om\cap B_{k+2})-E(\Om)
   \;\le\;\tfrac12\,\abs{\Om\setminus B_{k+1}}
   \Bigl(\sup_{\Om\setminus B_{k+1}}u_\Om
        +\sup_{\Om\setminus B_{k+1}}u_\Om^2\Bigr).
\end{equation}
%% JUSTIFICATION.  On B_{k+1}, \varphi_k=1 so (1-\varphi_k)=0 and
%% \nabla\varphi_k=0; thus both integrands vanish on B_{k+1} and on
%% \R^n\setminus B_{k+2} (where \varphi_k=0 is constant, gradient 0; and
%% u_\Om may be nonzero only inside \Om, but (1-\varphi_k)^2 there equals
%% 1 --- careful: outside B_{k+2}, 1-\varphi_k=1, NOT 0).  The precise
%% support is: \nabla\varphi_k is supported in the annulus
%% B_{k+2}\setminus B_{k+1}\subset\Om\setminus B_{k+1}; and (1-\varphi_k)^2
%% is supported in \R^n\setminus B_{k+1}.  Both correction integrals are
%% therefore over \Om\setminus B_{k+1}.  Bounding 0\le(1-\varphi_k)\le1
%% and |\nabla\varphi_k|\le1 and pulling out the sup of u_\Om (resp.
%% u_\Om^2) over that region gives
%%   \tfrac12\int_{\Om\setminus B_{k+1}}u_\Om+\tfrac12\int_{\Om\setminus
%%   B_{k+1}}u_\Om^2\le\tfrac12|\Om\setminus B_{k+1}|(\sup u_\Om+\sup
%%   u_\Om^2),
%% which is \eqref{eq:energy-comp-pre}.
By Lemma~\ref{lem:Linfty} applied with $R=k\ge1$ (so $R+1=k+1$ and the
tail region is $\Om\setminus B_{k+1}$),
$\sup_{\Om\setminus B_{k+1}}u_\Om\le C_\flat\abs{\Om\setminus B_k}^{1/n}
=C_\flat\,\omega_n^{1/n}b_k^{1/n}$, and likewise the square is bounded by
$C_\flat^2\omega_n^{2/n}b_k^{2/n}$. Also
$\abs{\Om\setminus B_{k+1}}=\omega_n b_{k+1}\le\omega_n b_k$. Plugging in,
and using $b_k\le1$ so that $b_k^{2/n}\le b_k^{1/n}$,
\begin{equation}\label{eq:energy-comp}
   E(\Om\cap B_{k+2})\;\le\;E(\Om)+C\,b_{k+1}\bigl(b_k^{1/n}+b_k^{2/n}\bigr)
   \;\le\;E(\Om)+C\,b_k^{\,1+1/n},
\end{equation}
for a dimensional constant $C=C(n)$ (we absorbed
$\tfrac12\omega_n(C_\flat\omega_n^{1/n}+C_\flat^2\omega_n^{2/n})$ and
$b_{k+1}\le b_k$ into $C$). This is the energy-comparison inequality.
%% JUSTIFICATION.  Lemma~\ref{lem:Linfty} with R=k (\ge1) gives
%% \sup_{\Om\setminus B_{k+1}}u_\Om\le C_\flat|\Om\setminus B_k|^{1/n};
%% and |\Om\setminus B_k|=\omega_n b_k so the bound is C_\flat\omega_n^{1/n}
%% b_k^{1/n}.  Squaring: \sup u_\Om^2=(\sup u_\Om)^2\le C_\flat^2
%% \omega_n^{2/n}b_k^{2/n}.  And |\Om\setminus B_{k+1}|=\omega_n b_{k+1}.
%% Substituting into \eqref{eq:energy-comp-pre} and pulling out the
%% dimensional prefactor:
%%   E(\Om\cap B_{k+2})-E(\Om)\le\tfrac12\omega_n b_{k+1}
%%     (C_\flat\omega_n^{1/n}b_k^{1/n}+C_\flat^2\omega_n^{2/n}b_k^{2/n})
%%     \le C b_{k+1}(b_k^{1/n}+b_k^{2/n}).
%% Since b_k\le1, b_k^{2/n}\le b_k^{1/n} (smaller positive base to a
%% larger... here 2/n vs 1/n: for n\ge2, 2/n\ge1/n, and base \le1 to a
%% LARGER exponent is SMALLER), so b_k^{1/n}+b_k^{2/n}\le2 b_k^{1/n}; and
%% b_{k+1}\le b_k.  Hence
%%   \le C\,b_k\cdot2 b_k^{1/n}=2C b_k^{1+1/n}, absorb 2 into C.
%% NOTE on n=2: replace 1/n by \kappa=1/4 throughout (Remark~\ref{rem:n2}),
%% giving the exponent 1+\kappa; nothing else changes.

\smallskip
\emph{Step 2: turning energy comparison into a tail recursion.}
By the Saint--Venant inequality \eqref{eq:saint-venant} applied to the
set $\Om\cap B_{k+2}$, whose volume is
$\abs{\Om\cap B_{k+2}}=\omega_n(1-b_{k+2})$,
\[
   E(B_1)\abs{B_1}^{-(n+2)/n}
   \le E(\Om\cap B_{k+2})\,\abs{\Om\cap B_{k+2}}^{-(n+2)/n},
\]
that is
\[
   E(B_1)\,\omega_n^{-(n+2)/n}\bigl(\omega_n(1-b_{k+2})\bigr)^{(n+2)/n}
   \le E(\Om\cap B_{k+2}),
\]
i.e.
\[
   E(B_1)\,(1-b_{k+2})^{(n+2)/n}\le E(\Om\cap B_{k+2}).
\]
%% JUSTIFICATION.  |\Om\cap B_{k+2}|=|\Om|-|\Om\setminus B_{k+2}|
%% =\omega_n-\omega_n b_{k+2}=\omega_n(1-b_{k+2}).  Saint--Venant
%% \eqref{eq:saint-venant} with the ball B_1:
%%   E(B_1)\omega_n^{-(n+2)/n}\le E(\Om\cap B_{k+2})|\Om\cap B_{k+2}|^{-(n+2)/n}.
%% Multiply both sides by |\Om\cap B_{k+2}|^{(n+2)/n}=
%% (\omega_n(1-b_{k+2}))^{(n+2)/n}=\omega_n^{(n+2)/n}(1-b_{k+2})^{(n+2)/n}.
%% The \omega_n^{(n+2)/n} cancels the \omega_n^{-(n+2)/n} on the left,
%% leaving E(B_1)(1-b_{k+2})^{(n+2)/n}\le E(\Om\cap B_{k+2}).
%% (Sign: E(B_1)<0 and (1-b_{k+2})^{(n+2)/n}\le1, so the left side is a
%% NEGATIVE number that is \ge E(B_1); the inequality direction is correct
%% because we multiplied by a positive quantity.)
Combine with \eqref{eq:energy-comp} and the definition
$E(\Om)=E(B_1)+\omega_n^{(n+2)/n}D(\Om)$ (which is
\eqref{eq:deficit-def} at $\abs{\Om}=\omega_n$):
\[
   E(B_1)(1-b_{k+2})^{(n+2)/n}
   \le E(\Om)+C\,b_k^{1+1/n}
   = E(B_1)+\omega_n^{(n+2)/n}D(\Om)+C\,b_k^{1+1/n}.
\]
%% JUSTIFICATION.  Chain the two inequalities just derived:
%%   E(B_1)(1-b_{k+2})^{(n+2)/n}\le E(\Om\cap B_{k+2})\le E(\Om)+C b_k^{1+1/n}
%% (the first from Saint--Venant, the second from \eqref{eq:energy-comp}).
%% Then substitute E(\Om)=E(B_1)+\omega_n^{(n+2)/n}D(\Om), the rearranged
%% definition of D at unit volume.
Rearranging, and recalling $E(B_1)<0$, divide by $-E(B_1)>0$:
\[
   1-(1-b_{k+2})^{(n+2)/n}
   \le\frac{\omega_n^{(n+2)/n}D(\Om)+C\,b_k^{1+1/n}}{-E(B_1)}
   =:\widehat C_1\bigl(D(\Om)+b_k^{1+1/n}\bigr),
\]
with $\widehat C_1=\widehat C_1(n)>0$
(here $\widehat C_1=\max\{\omega_n^{(n+2)/n},C\}/(-E(B_1))$).
%% JUSTIFICATION.  Subtract E(B_1) from both sides of the previous
%% display:
%%   E(B_1)[(1-b_{k+2})^{(n+2)/n}-1]\le\omega_n^{(n+2)/n}D(\Om)+C b_k^{1+1/n}.
%% The left side is E(B_1)\cdot(\text{negative})=positive\times... since
%% E(B_1)<0 and the bracket \le0, the product is \ge0; rewrite as
%%   (-E(B_1))[1-(1-b_{k+2})^{(n+2)/n}]\le\omega_n^{(n+2)/n}D+C b_k^{1+1/n}.
%% Divide by -E(B_1)>0 (preserves the inequality):
%%   1-(1-b_{k+2})^{(n+2)/n}\le[\omega_n^{(n+2)/n}D+C b_k^{1+1/n}]/(-E(B_1)).
%% Bound each coefficient by \max\{\omega_n^{(n+2)/n},C\}/(-E(B_1))=:
%% \widehat C_1, factoring out (D+b_k^{1+1/n}).  \widehat C_1>0 since
%% -E(B_1)>0.
To extract $b_{k+2}$ we use the elementary lower bound
\begin{equation}\label{eq:power-lb}
   1-(1-t)^{p}\;\ge\;t\qquad\text{for all }t\in[0,1],\ p\ge1,
\end{equation}
which holds because $r\mapsto(1-t)^{r}$ is non-increasing on $[0,\infty)$,
so $(1-t)^{p}\le(1-t)^{1}=1-t$ for $p=(n+2)/n\ge1$. Applying
\eqref{eq:power-lb} with $t=b_{k+2}\in[0,1]$ and $p=(n+2)/n$,
\begin{equation}\label{eq:bk2-bound}
   b_{k+2}\;\le\;1-(1-b_{k+2})^{(n+2)/n}
   \;\le\;\widehat C\,\bigl(D(\Om)+b_k^{\,1+1/n}\bigr),
   \qquad \widehat C:=\widehat C_1 .
\end{equation}
%% JUSTIFICATION (\eqref{eq:power-lb}).  For fixed t\in[0,1], the base
%% 1-t\in[0,1], and x\mapsto(1-t)^x is nonincreasing in x\ge0 (a base in
%% [0,1] to a larger power is smaller).  So for p\ge1, (1-t)^p\le(1-t)^1
%% =1-t, hence 1-(1-t)^p\ge1-(1-t)=t.  With t=b_{k+2}\in[0,1] (b's lie in
%% [0,1]) and p=(n+2)/n\ge1: b_{k+2}\le1-(1-b_{k+2})^{(n+2)/n}.  Chain
%% with the previous display to get b_{k+2}\le\widehat C(D+b_k^{1+1/n}).
%% This is the tail recursion: b two radii out \le const\times(deficit +
%% b here to a super-linear power).

\smallskip
\emph{Step 3: the geometric iteration to a bounded stopping index.}
Define the stopping index
\begin{equation}\label{eq:K-def}
   K:=\max\bigl\{k\ge1:\ b_k\ge2\widehat C\,D(\Om)\bigr\},
\end{equation}
with the convention $K:=0$ if $b_1<2\widehat C D(\Om)$ (then the tail is
already controlled at the first radius and Step~5 applies with $K=0$).
The set in \eqref{eq:K-def} is finite because $b_k\to0$ as $k\to\infty$
(the tail mass of a finite-measure set vanishes), so $K$ is well defined.
%% JUSTIFICATION.  b_k=|\Om\setminus B_k|/\omega_n\to0 since
%% |\Om\setminus B_k|\downarrow|\Om\setminus\bigcup_k B_k|=|\Om\setminus
%% \R^n|=0 (or directly: |\Om|<\infty forces the tails to vanish).  Hence
%% only finitely many k have b_k\ge2\widehat C D(\Om)>0 (note D(\Om)>0 in
%% the nontrivial case; if D(\Om)=0 then \Om=B_1 a.e.\ by rigidity and
%% there is nothing to prove).  The max of a finite nonempty set exists;
%% if the set is empty, K=0 by convention.
We prove
\begin{equation}\label{eq:K-bound}
   K\;\le\;K(n)
\end{equation}
for a dimensional $K(n)$, provided $\delta(n)$ is small enough. Assume
$K\ge3$ (otherwise \eqref{eq:K-bound} is trivial).
%% JUSTIFICATION.  If K\le2 then K\le K(n) for any K(n)\ge2, so we may
%% assume K\ge3 to run the odd-index iteration (which needs at least the
%% indices 1 and 3 in range).

\emph{Self-improving contraction.} For every $k$ with $k+2\le K$ we have,
by definition of $K$, that $b_{k+2}\ge2\widehat C D(\Om)$, i.e.\
$\widehat C D(\Om)\le\tfrac12 b_{k+2}$. Feeding this into
\eqref{eq:bk2-bound},
\[
   b_{k+2}\le\widehat C D(\Om)+\widehat C b_k^{1+1/n}
   \le\tfrac12 b_{k+2}+\widehat C b_k^{1+1/n},
\]
and absorbing $\tfrac12 b_{k+2}$ on the left,
\begin{equation}\label{eq:bk2-self}
   b_{k+2}\;\le\;2\widehat C\,b_k^{\,1+1/n}
   \qquad\text{whenever }k+2\le K .
\end{equation}
%% JUSTIFICATION.  k+2\le K and K is the LARGEST index with b\ge2\widehat
%% C D, so b_{k+2}\ge2\widehat C D(\Om), i.e. \widehat C D(\Om)\le
%% \tfrac12 b_{k+2}.  \eqref{eq:bk2-bound} says b_{k+2}\le\widehat C D+
%% \widehat C b_k^{1+1/n}; bound the first summand by \tfrac12 b_{k+2}:
%%   b_{k+2}\le\tfrac12 b_{k+2}+\widehat C b_k^{1+1/n}.
%% Subtract \tfrac12 b_{k+2}: \tfrac12 b_{k+2}\le\widehat C b_k^{1+1/n},
%% i.e. b_{k+2}\le2\widehat C b_k^{1+1/n}.  This is the self-improvement:
%% as long as we are below the stopping index, the deficit term is
%% subdominant and the recursion is purely power-contractive.

\emph{Smallness from the seed.} By Step~0, $b_1\le(C_\dagger D(\Om))^{1/4}
\le(C_\dagger\delta(n))^{1/4}$. Choose $\delta(n)$ small enough that
\begin{equation}\label{eq:seed-smallness}
   b_1\;\le\;(C_\dagger\,\delta(n))^{1/4}\;\le\;(2\widehat C)^{-2n}
   \;=:\;\eta_n\in(0,1).
\end{equation}
Then $2\widehat C\,b_k^{1/(2n)}\le2\widehat C\,\eta_n^{1/(2n)}=1$ for all
$k\ge1$ (as $b_k\le b_1\le\eta_n$), so \eqref{eq:bk2-self} self-improves
to
\begin{equation}\label{eq:bk2-power}
   b_{k+2}\;=\;\bigl(2\widehat C\,b_k^{1/(2n)}\bigr)\,b_k^{\,1+1/(2n)}
   \;\le\;b_k^{\,1+1/(2n)}\qquad\text{whenever }k+2\le K .
\end{equation}
%% JUSTIFICATION.  Choose \delta(n) so that (C_\dagger\delta(n))^{1/4}\le
%% (2\widehat C)^{-2n}=:\eta_n (possible since the left side \to0 as
%% \delta\to0; and \eta_n\in(0,1) because 2\widehat C>1, say, or at any
%% rate (2\widehat C)^{-2n}\in(0,1) for \widehat C\ge1, which we may
%% assume by enlarging it).  Then for all k, b_k\le b_1\le\eta_n=
%% (2\widehat C)^{-2n}, so b_k^{1/(2n)}\le(2\widehat C)^{-1}, hence
%% 2\widehat C b_k^{1/(2n)}\le1.  Now SPLIT the exponent in
%% \eqref{eq:bk2-self}: write b_k^{1+1/n}=b_k^{1/(2n)}\cdot b_k^{1+1/(2n)}
%% (check exponents: 1/(2n)+1+1/(2n)=1+1/n).  Thus
%%   b_{k+2}\le2\widehat C b_k^{1+1/n}=(2\widehat C b_k^{1/(2n)})
%%     b_k^{1+1/(2n)}\le1\cdot b_k^{1+1/(2n)}=b_k^{1+1/(2n)}.
%% The constant 2\widehat C has been completely absorbed by the smallness
%% of b_k; the recursion is now a clean power map with no constant.

\emph{Iterating along odd indices.} Apply \eqref{eq:bk2-power} repeatedly,
starting from $k=1$, as long as the index stays $\le K$. Writing the odd
indices $1,3,5,\dots$ and setting $\gamma:=1+\tfrac1{2n}>1$, induction
gives, for every $j\ge0$ with $2j+1\le K$,
\begin{equation}\label{eq:bodd-decay}
   b_{2j+1}\;\le\;b_1^{\,\gamma^{\,j}} .
\end{equation}
Indeed \eqref{eq:bodd-decay} holds at $j=0$, and if it holds at $j$ with
$2(j+1)+1=2j+3\le K$, then \eqref{eq:bk2-power} (at $k=2j+1$, legitimate
since $2j+3\le K$) gives
$b_{2j+3}\le b_{2j+1}^{\gamma}\le\bigl(b_1^{\gamma^{j}}\bigr)^{\gamma}
=b_1^{\gamma^{j+1}}$.
%% JUSTIFICATION (induction).  Base j=0: b_1\le b_1^{\gamma^0}=b_1^1,
%% equality.  Step: \eqref{eq:bk2-power} at k=2j+1 requires k+2=2j+3\le K,
%% which is the hypothesis; it gives b_{2j+3}=b_{(2j+1)+2}\le b_{2j+1}^{\gamma}.
%% By the inductive hypothesis b_{2j+1}\le b_1^{\gamma^j}, and raising to
%% the power \gamma>0 (monotone, base\le1) gives b_{2j+1}^{\gamma}\le
%% (b_1^{\gamma^j})^{\gamma}=b_1^{\gamma^{j+1}}.  Chain.  Note the DOUBLY
%% EXPONENTIAL decay: the exponent on b_1<1 grows like \gamma^j, so b's
%% crash extremely fast --- this is what bounds K.

\emph{Bounding $K$.} Let $j_K:=\lfloor(K-1)/2\rfloor$, the largest $j$
with $2j+1\le K$. By the maximality of $K$ and \eqref{eq:bodd-decay},
\begin{equation}\label{eq:K-pinch}
   2\widehat C\,D(\Om)\;\le\;b_{2j_K+1}\;\le\;b_1^{\,\gamma^{\,j_K}}
   \;\le\;\bigl(C_\dagger D(\Om)\bigr)^{\gamma^{\,j_K}/4},
\end{equation}
using $b_1\le(C_\dagger D(\Om))^{1/4}$.
%% JUSTIFICATION.  2j_K+1\le K, so by the DEFINITION of K (every index
%% up to K has b\ge2\widehat C D) we have b_{2j_K+1}\ge2\widehat C D(\Om).
%% By \eqref{eq:bodd-decay} at j=j_K, b_{2j_K+1}\le b_1^{\gamma^{j_K}}.
%% And b_1\le(C_\dagger D)^{1/4}, so b_1^{\gamma^{j_K}}\le
%% ((C_\dagger D)^{1/4})^{\gamma^{j_K}}=(C_\dagger D)^{\gamma^{j_K}/4}.
%% Chain the three.  This "pinch" traps 2\widehat C D between the deficit
%% scale and a very high power of the deficit; for D small that forces
%% the exponent \gamma^{j_K} to be bounded.
Take logarithms (recall
$D(\Om)\le\delta(n)<1$, so $\log D(\Om)<0$); writing $L:=-\log D(\Om)>0$,
\eqref{eq:K-pinch} reads
\[
   \log(2\widehat C)-L
   \;\le\;\frac{\gamma^{\,j_K}}{4}\bigl(\log C_\dagger-L\bigr).
\]
%% JUSTIFICATION (taking logs; sign handling).  Take \log of the outer
%% inequality 2\widehat C D\le(C_\dagger D)^{\gamma^{j_K}/4} in
%% \eqref{eq:K-pinch}.  \log is increasing, so the inequality DIRECTION
%% is preserved:
%%   \log(2\widehat C)+\log D\le(\gamma^{j_K}/4)(\log C_\dagger+\log D).
%% Now \log D=-L (L>0 since 0<D<1).  Substitute:
%%   \log(2\widehat C)-L\le(\gamma^{j_K}/4)(\log C_\dagger-L).
%% This is the displayed inequality.  Crucially we did NOT divide by a
%% negative number yet; both sides may be negative, but the log step is
%% safe because log is monotone.
For $\delta(n)$ small, $L=-\log D(\Om)\ge-\log\delta(n)$ is large; in
particular we may take $\delta(n)$ small enough that
$L\ge2\max\{\abs{\log(2\widehat C)},\abs{\log C_\dagger}\}$, whence
$\log(2\widehat C)-L\ge-\tfrac32 L$ and $\log C_\dagger-L\le-\tfrac12 L$.
%% JUSTIFICATION.  D\le\delta(n) gives L=-\log D\ge-\log\delta(n), which
%% \to+\infty as \delta\to0; so we may force L\ge2M_0 where
%% M_0:=\max\{|\log(2\widehat C)|,|\log C_\dagger|\}.  Then:
%%  \log(2\widehat C)\ge-|\log(2\widehat C)|\ge-M_0\ge-\tfrac12 L, so
%%    \log(2\widehat C)-L\ge-\tfrac12 L-L=-\tfrac32 L;
%%  \log C_\dagger\le|\log C_\dagger|\le M_0\le\tfrac12 L, so
%%    \log C_\dagger-L\le\tfrac12 L-L=-\tfrac12 L.
%% Both bounds are sign-clean: the LHS of the log-inequality is \ge-\tfrac32 L
%% and the bracket on the RHS is \le-\tfrac12 L (negative).
The displayed inequality then forces
$-\tfrac32 L\le-\tfrac{\gamma^{\,j_K}}{8}L$, i.e.\
$\gamma^{\,j_K}\le12$.
%% JUSTIFICATION (the sign handling on dividing).  From the log-inequality
%%   \log(2\widehat C)-L\le(\gamma^{j_K}/4)(\log C_\dagger-L)
%% bound the LHS below and the bracket above:
%%   -\tfrac32 L\le\log(2\widehat C)-L\le(\gamma^{j_K}/4)(\log C_\dagger-L)
%%             \le(\gamma^{j_K}/4)(-\tfrac12 L)=-\tfrac{\gamma^{j_K}}{8}L.
%% The last \le uses that the bracket (\log C_\dagger-L) is NEGATIVE
%% (\le-\tfrac12 L<0) and \gamma^{j_K}/4>0, so multiplying the bound
%% (\log C_\dagger-L)\le-\tfrac12 L by the POSITIVE \gamma^{j_K}/4
%% preserves direction.  Now we have -\tfrac32 L\le-\tfrac{\gamma^{j_K}}{8}L.
%% Divide by -L<0, which REVERSES the inequality:
%%   \tfrac32\ge\tfrac{\gamma^{j_K}}{8}, i.e. \gamma^{j_K}\le12.
%% (This is the one place dividing by a negative number flips the sign;
%% we tracked it explicitly.)
Since $\gamma=1+\tfrac1{2n}>1$ this bounds
$j_K\le\log12/\log\gamma=:j_*(n)$, hence
$K\le2j_K+2\le2j_*(n)+2=:K(n)$, a dimensional constant. This proves
\eqref{eq:K-bound}.
%% JUSTIFICATION.  \gamma^{j_K}\le12 with \gamma>1: take logs (both sides
%% positive), j_K\log\gamma\le\log12, so j_K\le\log12/\log\gamma=:j_*(n)
%% (depends on n through \gamma=1+1/(2n)).  Since j_K=\lfloor(K-1)/2\rfloor
%% \ge(K-1)/2-1=(K-3)/2, we get (K-3)/2\le j_K\le j_*(n)?  More simply
%% K\le2j_K+2 (as 2j_K+1\le K\le2(j_K+1)+1-1=2j_K+2 from maximality of
%% j_K), so K\le2j_*(n)+2=:K(n).  Dimensional.

\emph{Step 4: tail at the stopping radius.} By the maximality
\eqref{eq:K-def}, $b_{K+1}<2\widehat C D(\Om)$; we record this for the next
step.
%% JUSTIFICATION.  K is the LARGEST k with b_k\ge2\widehat C D, so the
%% next index K+1 fails the condition: b_{K+1}<2\widehat C D(\Om).  (If
%% K=0 by convention, this is exactly b_1<2\widehat C D, also fine.)

\smallskip
\emph{Step 5: the truncated set and its volume.}
By the definition \eqref{eq:K-def} of $K$ and $K\le K(n)$,
$b_{K+1}<2\widehat C D(\Om)$, hence
\begin{equation}\label{eq:tail-controlled}
   \abs{\Om\setminus B_{K+1}}=\omega_n b_{K+1}<2\widehat C\,\omega_n\,D(\Om).
\end{equation}
In fact we cut one radius further out to have room for the cut-off: by
monotonicity $b_{K+3}\le b_{K+1}<2\widehat C D(\Om)$, so
\begin{equation}\label{eq:vol-cut}
   \abs{\Om\cap B_{K+3}}=\omega_n(1-b_{K+3})\ge\omega_n(1-2\widehat C D(\Om)),
   \qquad
   \abs{\Om\setminus B_{K+3}}<2\widehat C\,\omega_n D(\Om).
\end{equation}
%% JUSTIFICATION.  \eqref{eq:tail-controlled}: multiply b_{K+1}<2\widehat
%% C D by \omega_n and use |\Om\setminus B_{K+1}|=\omega_n b_{K+1}.
%% \eqref{eq:vol-cut}: b_{K+3}\le b_{K+1} (monotonicity, K+3>K+1)<2\widehat
%% C D; then |\Om\cap B_{K+3}|=\omega_n(1-b_{K+3})\ge\omega_n(1-2\widehat C D)
%% and |\Om\setminus B_{K+3}|=\omega_n b_{K+3}<2\widehat C\omega_n D.
%% We cut at K+3 rather than K+1 so that the energy-comparison cut-off
%% \varphi_{K+1} (which lives on the annulus B_{K+3}\setminus B_{K+2}) has
%% room; this is the "+2" built into Step 1's cut-off.
Define
\[
   r:=\Bigl(\frac{\abs{\Om\cap B_{K+3}}}{\omega_n}\Bigr)^{1/n}
   =(1-b_{K+3})^{1/n}\in\bigl[(1-2\widehat C D(\Om))^{1/n},\,1\bigr],
   \qquad
   \Otilde:=r^{-1}\,(\Om\cap B_{K+3}).
\]
Then $\abs{\Otilde}=r^{-n}\abs{\Om\cap B_{K+3}}=\omega_n$, which is
\eqref{eq:trunc-vol}.
%% JUSTIFICATION.  r=(|\Om\cap B_{K+3}|/\omega_n)^{1/n}=(1-b_{K+3})^{1/n};
%% since 0\le b_{K+3}<2\widehat C D\le2\widehat C\delta, r lies in
%% [(1-2\widehat C D)^{1/n},1].  Dilation by r^{-1} scales volume by
%% r^{-n}: |\Otilde|=r^{-n}|\Om\cap B_{K+3}|=r^{-n}\cdot\omega_n r^n
%% =\omega_n.  This rescaling restores unit volume after the cut removed a
%% sliver; r is chosen exactly to do that.
Since $\Om\cap B_{K+3}\subset B_{K+3}$ and
$r\ge(1-2\widehat C\delta(n))^{1/n}\ge\tfrac12$ for $\delta(n)$ small,
\[
   \diam(\Otilde)=r^{-1}\diam(\Om\cap B_{K+3})\le r^{-1}\cdot2(K+3)
   \le4\bigl(K(n)+3\bigr)=:d(n),
\]
giving \eqref{eq:trunc-diam}.
%% JUSTIFICATION.  diam scales linearly: diam(r^{-1}S)=r^{-1}diam(S).
%% \Om\cap B_{K+3}\subset B_{K+3} has diam\le2(K+3).  r\ge(1-2\widehat C
%% \delta)^{1/n}\ge1/2 for \delta small (since (1-x)^{1/n}\to1 as x\to0;
%% pick \delta with 2\widehat C\delta small enough that the n-th root is
%% \ge1/2, e.g. 1-2\widehat C\delta\ge2^{-n}).  So r^{-1}\le2 and
%%   diam(\Otilde)\le2\cdot2(K+3)=4(K+3)\le4(K(n)+3)=:d(n),
%% using K\le K(n).  d(n) dimensional.
We record for later the linearised bound
\begin{equation}\label{eq:r-near-1}
   0\le1-r^n=b_{K+3}<2\widehat C D(\Om),
   \qquad
   0\le1-r\le1-r^n<2\widehat C D(\Om),
\end{equation}
the middle inequality because $r\le1$ implies $1-r\le1-r^n$.
%% JUSTIFICATION.  1-r^n=1-(1-b_{K+3})=b_{K+3}\in[0,2\widehat C D).  For
%% the second chain: r\in(0,1], and for such r and n\ge1, r^n\le r, so
%% 1-r\le1-r^n.  Both are <2\widehat C D.  These are the "r is close to 1"
%% bounds used in Steps 6 and 7.

\smallskip
\emph{Step 6: deficit of the surrogate.}
Apply the energy comparison \eqref{eq:energy-comp} at $k=K+1$ (so the
cut radius is $k+2=K+3$):
\[
   E(\Om\cap B_{K+3})\le E(\Om)+C\,b_{K+1}^{\,1+1/n}
   \le E(\Om)+C\,(2\widehat C D(\Om))^{1+1/n}
   \le E(\Om)+C'\,D(\Om),
\]
where in the last step we used $D(\Om)\le\delta(n)\le1$ so that
$D(\Om)^{1+1/n}\le D(\Om)$, and absorbed constants into
$C'=C'(n)$.
%% JUSTIFICATION.  \eqref{eq:energy-comp} at k=K+1 reads E(\Om\cap B_{K+3})
%% \le E(\Om)+C b_{K+1}^{1+1/n}.  Use b_{K+1}<2\widehat C D
%% (\eqref{eq:tail-controlled}) and t\mapsto t^{1+1/n} increasing:
%% b_{K+1}^{1+1/n}\le(2\widehat C D)^{1+1/n}=(2\widehat C)^{1+1/n}D^{1+1/n}.
%% Since 0\le D\le1 and 1+1/n\ge1, D^{1+1/n}\le D.  So the correction is
%% \le C(2\widehat C)^{1+1/n}D=:C'D, C'=C'(n).
Since $E$ scales as
$E(t\Omega')=t^{n+2}E(\Omega')$ and $\Otilde=r^{-1}(\Om\cap B_{K+3})$,
\[
   E(\Otilde)=r^{-(n+2)}E(\Om\cap B_{K+3})
   \le r^{-(n+2)}\bigl(E(\Om)+C'D(\Om)\bigr).
\]
%% JUSTIFICATION.  \Otilde=r^{-1}(\Om\cap B_{K+3}), so by the scaling law
%% (proved in \S\ref{ssec:notation}) E(\Otilde)=(r^{-1})^{n+2}
%% E(\Om\cap B_{K+3})=r^{-(n+2)}E(\Om\cap B_{K+3}).  Then substitute the
%% bound on E(\Om\cap B_{K+3}); r^{-(n+2)}>0 preserves the inequality.
Because $\abs{\Otilde}=\omega_n$,
$D(\Otilde)=\omega_n^{-(n+2)/n}(E(\Otilde)-E(B_1))$. Using
$E(\Om)=E(B_1)+\omega_n^{(n+2)/n}D(\Om)$ and $r\le1$ (so
$r^{-(n+2)}\ge1$), with $E(B_1)<0$,
\begin{align*}
   E(\Otilde)-E(B_1)
   &\le r^{-(n+2)}\bigl(E(B_1)+\omega_n^{(n+2)/n}D(\Om)+C'D(\Om)\bigr)-E(B_1)\\
   &=\bigl(r^{-(n+2)}-1\bigr)E(B_1)
     +r^{-(n+2)}\bigl(\omega_n^{(n+2)/n}+C'\bigr)D(\Om).
\end{align*}
%% JUSTIFICATION.  Substitute E(\Om)=E(B_1)+\omega_n^{(n+2)/n}D into the
%% bound on E(\Otilde) and subtract E(B_1):
%%   E(\Otilde)-E(B_1)\le r^{-(n+2)}(E(B_1)+\omega_n^{(n+2)/n}D+C'D)-E(B_1).
%% Group the E(B_1) terms: r^{-(n+2)}E(B_1)-E(B_1)=(r^{-(n+2)}-1)E(B_1);
%% and the D terms: r^{-(n+2)}(\omega_n^{(n+2)/n}+C')D.  This is the
%% second line.
For the first term, $r^{-(n+2)}-1\ge0$ and $E(B_1)<0$, so
$(r^{-(n+2)}-1)E(B_1)\le0$ and we may bound it by
$(r^{-(n+2)}-1)|E(B_1)|$.
%% JUSTIFICATION.  r\le1\Rightarrow r^{-(n+2)}\ge1\Rightarrow
%% r^{-(n+2)}-1\ge0; times E(B_1)<0 gives a NONPOSITIVE number, which is
%% \le its absolute value (r^{-(n+2)}-1)|E(B_1)|.  (A nonpositive number
%% is \le any nonnegative number; we use this nonnegative surrogate as a
%% clean upper bound.)
Quantitatively, by \eqref{eq:r-near-1},
$r^{n}\ge1-2\widehat C D(\Om)$, so
$r^{-(n+2)}=r^{-n}\cdot r^{-2}\le(1-2\widehat C D(\Om))^{-(n+2)/n}
\le1+C''D(\Om)$ for $\delta(n)$ small (Bernoulli/Taylor bound, since
$2\widehat C D(\Om)\le2\widehat C\delta(n)\le\tfrac12$). Hence
$r^{-(n+2)}-1\le C''D(\Om)$ and $r^{-(n+2)}\le1+C''D(\Om)\le2$, giving
\[
   E(\Otilde)-E(B_1)
   \le C''D(\Om)\cdot\abs{E(B_1)}+2\bigl(\omega_n^{(n+2)/n}+C'\bigr)D(\Om)
   \le C_2'\,D(\Om),
\]
with $C_2'=C_2'(n)$.
%% JUSTIFICATION (the Bernoulli/Taylor bound).  From \eqref{eq:r-near-1},
%% r^n=1-b_{K+3}\ge1-2\widehat C D.  So r^{-n}=(r^n)^{-1}\le(1-2\widehat C
%% D)^{-1}, and r^{-(n+2)}=(r^n)^{-(n+2)/n}\le(1-2\widehat C D)^{-(n+2)/n}.
%% Set x:=2\widehat C D\in[0,1/2] (since 2\widehat C D\le2\widehat C\delta
%% \le1/2 for \delta small).  For x\in[0,1/2] and exponent q:=(n+2)/n\le3,
%% (1-x)^{-q}\le1+C''x with C''=C''(n) (a one-sided Taylor/Bernoulli bound:
%% e.g. (1-x)^{-q}\le1+2q x for x\le1/2, so C''=2q\cdot2\widehat C works).
%% Thus r^{-(n+2)}\le1+C''D, hence r^{-(n+2)}-1\le C''D and r^{-(n+2)}\le
%% 1+C''D\le1+C''\delta\le2 for \delta small.  Plug both into the two
%% terms: first \le C''D|E(B_1)|, second \le2(\omega_n^{(n+2)/n}+C')D.
%% Sum and absorb into C_2'(n).
Dividing by $\omega_n^{(n+2)/n}$,
\begin{equation}\label{eq:D-tilde}
   D(\Otilde)=\omega_n^{-(n+2)/n}\bigl(E(\Otilde)-E(B_1)\bigr)
   \le\omega_n^{-(n+2)/n}C_2'\,D(\Om)=:C_3\,D(\Om),
\end{equation}
which is \eqref{eq:trunc-D} with $C_3=C_3(n)$.
%% JUSTIFICATION.  D(\Otilde)=\omega_n^{-(n+2)/n}(E(\Otilde)-E(B_1))
%% (unit volume); apply the bound E(\Otilde)-E(B_1)\le C_2'D and set
%% C_3=\omega_n^{-(n+2)/n}C_2'(n).

\smallskip
\emph{Step 7: asymmetry loss.}
Let $B_1(x_0)$ realise the asymmetry of $\Otilde$, so
$\Fr(\Otilde)=\abs{\Otilde\,\triangle\,B_1(x_0)}/\omega_n$. We compare the
asymmetry of $\Om$ to that of $\Otilde$ through the chain of inclusions
relating $\Om$, $\Om\cap B_{K+3}$ and $\Otilde=r^{-1}(\Om\cap B_{K+3})$.
For any test ball $B_1(y)$,
\[
   \abs{\Om\,\triangle\,B_1(y)}
   \le\abs{\Om\,\triangle\,(\Om\cap B_{K+3})}
     +\abs{(\Om\cap B_{K+3})\,\triangle\,B_1(y)}.
\]
%% JUSTIFICATION (triangle inequality for symmetric difference).  The map
%% A\mapsto\mathbf 1_A is an isometry of (sets, \triangle) into L^1, with
%% |A\triangle B|=\|\mathbf 1_A-\mathbf 1_B\|_{L^1}.  Hence \triangle
%% obeys the triangle inequality: for any A,M,B,
%%   |A\triangle B|\le|A\triangle M|+|M\triangle B|.
%% Take A=\Om, M=\Om\cap B_{K+3}, B=B_1(y).  This is the "asymmetry-loss
%% triangle inequality" the brief refers to.
The first term is $\abs{\Om\setminus B_{K+3}}<2\widehat C\omega_n D(\Om)$
by \eqref{eq:vol-cut}.
%% JUSTIFICATION.  \Om\triangle(\Om\cap B_{K+3})=\Om\setminus(\Om\cap
%% B_{K+3})=\Om\setminus B_{K+3} (since \Om\cap B_{K+3}\subset\Om, the
%% symmetric difference is just what \Om loses, i.e. \Om\setminus B_{K+3}).
%% Its measure is <2\widehat C\omega_n D by \eqref{eq:vol-cut}.
For the second, scale by $r$: writing
$E_r:=\Om\cap B_{K+3}=r\Otilde$ and $B_1(y)=r\,B_{1/r}(y/r)$,
\[
   \abs{(r\Otilde)\,\triangle\,(r\,B_{1/r}(y/r))}
   =r^n\,\abs{\Otilde\,\triangle\,B_{1/r}(y/r)}
   \le r^n\Bigl(\abs{\Otilde\,\triangle\,B_1(y/r)}
              +\abs{B_1(y/r)\,\triangle\,B_{1/r}(y/r)}\Bigr).
\]
%% JUSTIFICATION.  \Om\cap B_{K+3}=r\Otilde by definition of \Otilde
%% (\Otilde=r^{-1}(\Om\cap B_{K+3})).  A unit ball B_1(y)=\{|x-y|<1\}
%% equals r\{|z-y/r|<1/r\}=r B_{1/r}(y/r) (write x=rz).  Dilation by r
%% scales |\cdot\triangle\cdot| by r^n:
%%   |(rA)\triangle(rB)|=r^n|A\triangle B|.
%% So |(r\Otilde)\triangle(r B_{1/r}(y/r))|=r^n|\Otilde\triangle
%% B_{1/r}(y/r)|.  Then the triangle inequality inserts the intermediate
%% set B_1(y/r):
%%   |\Otilde\triangle B_{1/r}(y/r)|\le|\Otilde\triangle B_1(y/r)|
%%     +|B_1(y/r)\triangle B_{1/r}(y/r)|.
The last term is the volume difference of concentric balls of radii $1$
and $1/r\ge1$:
$\abs{B_1\,\triangle\,B_{1/r}}=\omega_n(r^{-n}-1)$. Hence
\[
   r^n\abs{B_1(y/r)\,\triangle\,B_{1/r}(y/r)}
   =r^n\,\omega_n(r^{-n}-1)=\omega_n(1-r^n)
   <2\widehat C\,\omega_n D(\Om),
\]
using \eqref{eq:r-near-1}.
%% JUSTIFICATION.  Two concentric balls of radii 1\le1/r: B_1\subset
%% B_{1/r}, so B_1\triangle B_{1/r}=B_{1/r}\setminus B_1, of volume
%% \omega_n((1/r)^n-1^n)=\omega_n(r^{-n}-1).  Multiply by r^n:
%%   r^n\omega_n(r^{-n}-1)=\omega_n(1-r^n).
%% By \eqref{eq:r-near-1}, 1-r^n<2\widehat C D, so this is <2\widehat C
%% \omega_n D.  (Translation by y/r does not change the volume.)
Choosing $y:=r x_0$ so that $y/r=x_0$ is the
optimal centre for $\Otilde$, the middle term becomes
$r^n\abs{\Otilde\,\triangle\,B_1(x_0)}=r^n\omega_n\Fr(\Otilde)
\le\omega_n\Fr(\Otilde)$ (as $r\le1$). Collecting,
\[
   \abs{\Om\,\triangle\,B_1(rx_0)}
   \le2\widehat C\omega_n D(\Om)+\omega_n\Fr(\Otilde)+2\widehat C\omega_n D(\Om),
\]
%% JUSTIFICATION.  Pick the test ball with y/r=x_0, i.e. y=rx_0; then the
%% middle term is r^n|\Otilde\triangle B_1(x_0)|=r^n\omega_n\Fr(\Otilde)
%% (definition of \Fr(\Otilde) with the OPTIMAL ball B_1(x_0)), and
%% r^n\le1 gives \le\omega_n\Fr(\Otilde).  Now assemble the three pieces:
%%   |\Om\triangle B_1(rx_0)|\le \underbrace{2\widehat C\omega_n D}_{1st}
%%     +\underbrace{\omega_n\Fr(\Otilde)}_{middle, scaled}
%%     +\underbrace{2\widehat C\omega_n D}_{conc. balls},
%% where "1st" is the cut loss and the third is the concentric-ball loss.
and dividing by $\omega_n$ and taking the infimum over translates (which
only decreases the left side),
\begin{equation}\label{eq:trunc-A-pre}
   \Fr(\Om)\le\frac{\abs{\Om\,\triangle\,B_1(rx_0)}}{\omega_n}
   \le\Fr(\Otilde)+4\widehat C\,D(\Om).
\end{equation}
This is \eqref{eq:trunc-A} with the constant $4\widehat C$.
%% JUSTIFICATION.  Divide the collected bound by \omega_n:
%%   |\Om\triangle B_1(rx_0)|/\omega_n\le\Fr(\Otilde)+4\widehat C D.
%% \Fr(\Om)=\inf_y|\Om\triangle B_1(y)|/\omega_n\le the value at the
%% particular y=rx_0, which is the displayed right side.  Hence
%% \Fr(\Om)\le\Fr(\Otilde)+4\widehat C D.  (The two cut/concentric losses
%% 2\widehat C D each summed to 4\widehat C D.)
Taking $C_\sharp:=\max\{C_3,4\widehat C\}=C_\sharp(n)$ proves
\eqref{eq:trunc-D}--\eqref{eq:trunc-A} simultaneously, and we already have
\eqref{eq:trunc-vol}--\eqref{eq:trunc-diam}. This completes the proof.
%% JUSTIFICATION.  \eqref{eq:trunc-D} holds with C_3, \eqref{eq:trunc-A}
%% with 4\widehat C; taking C_\sharp=\max\{C_3,4\widehat C\} makes both
%% read with the single constant C_\sharp.  Volume and diameter were
%% \eqref{eq:trunc-vol}, \eqref{eq:trunc-diam}.  All four hold.
\end{proof}

\begin{remark}[Translation into a fixed ball]\label{rem:translate}
The diameter bound \eqref{eq:trunc-diam} is translation invariant, so
after a translation we may assume $\Otilde\subset B_{d(n)}$. Both
$\Fr$ and $D$ are translation invariant, so
\eqref{eq:trunc-vol}--\eqref{eq:trunc-A} are preserved.
%% JUSTIFICATION.  A set of diameter \le d(n) fits inside a ball of radius
%% d(n) centred at any of its points; translate that point to 0, so
%% \Otilde\subset B_{d(n)}.  \Fr, D, |\cdot|, diam are all translation
%% invariant, so the four conclusions are unaffected.
\end{remark}

\subsection{The far branch}\label{ssec:far}

When the deficit is not small the inequality is trivial: the asymmetry is
bounded ($\Fr<2$) while the deficit is bounded below, so a quadratic
bound holds with a crude constant. No truncation is needed.

\begin{lemma}[Large-deficit quadratic bound]\label{lem:far}
For every open $\Om$ with $\abs{\Om}=\omega_n$ and
$D(\Om)\ge\delta(n)/4$,
\begin{equation}\label{eq:far}
   E(\Om)-E(B_1)\;\ge\;\omega_n^{(n+2)/n}\,\frac{\delta(n)}{16}\,\Fr(\Om)^2 .
\end{equation}
\end{lemma}

\begin{proof}
Since $\abs{\Om}=\omega_n$ gives $\Fr(\Om)<2$, we have
$\Fr(\Om)^2/4<1$. Hence
\[
   D(\Om)\ge\frac{\delta(n)}{4}
   =\frac{\delta(n)}{4}\cdot1
   \ge\frac{\delta(n)}{4}\cdot\frac{\Fr(\Om)^2}{4}
   =\frac{\delta(n)}{16}\Fr(\Om)^2 .
\]
%% JUSTIFICATION.  \Fr(\Om)<2 (asymmetry range at unit volume), so
%% \Fr(\Om)^2<4, i.e. \Fr(\Om)^2/4<1.  Start from D\ge\delta/4 and
%% multiply the "1" by the smaller quantity \Fr^2/4<1:
%%   D\ge\delta/4\ge(\delta/4)(\Fr^2/4)=(\delta/16)\Fr^2.
Multiplying by $\omega_n^{(n+2)/n}$ and using
$E(\Om)-E(B_1)=\omega_n^{(n+2)/n}D(\Om)$ gives \eqref{eq:far}.
%% JUSTIFICATION.  Multiply D\ge(\delta/16)\Fr^2 by \omega_n^{(n+2)/n}>0
%% and use the unit-volume identity E(\Om)-E(B_1)=\omega_n^{(n+2)/n}D.
\end{proof}

\subsection{The near branch}\label{ssec:near}

Fix the dimensional working radius
\begin{equation}\label{eq:Rstar}
   R_*(n):=\max\{d(n),2\},
\end{equation}
so $R_*\ge d(n)$ (the surrogates of Lemma~\ref{lem:truncation} lie in
$B_{R_*}$ after translation, by Remark~\ref{rem:translate}) and
$R_*\ge2$ (the hypothesis of Theorem~\ref{thm:SV-bounded}).
%% JUSTIFICATION.  We need a single radius that (a) contains every
%% surrogate \Otilde (so Theorem~\ref{thm:SV-bounded} applies to it) and
%% (b) meets the standing hypothesis R\ge2 of that theorem.  Taking the
%% max of d(n) and 2 secures both; R_*(n) is dimensional since d(n) is.

\begin{lemma}[Small-deficit transfer]\label{lem:near}
Set
\begin{equation}\label{eq:csv-tilde}
   \widetilde c_{\mathrm{SV}}(n)
   :=\omega_n^{-(n+2)/n}\,c_{\mathrm{SV}}\bigl(n,R_*(n)\bigr)>0 .
\end{equation}
For every open $\Om$ with $\abs{\Om}=\omega_n$ and
$D(\Om)\le\delta(n)$,
\begin{equation}\label{eq:near-Fr}
   \Fr(\Om)^2\;\le\;
   \Bigl(\frac{4\,C_\sharp(n)}{\widetilde c_{\mathrm{SV}}(n)}
        +4\,C_\sharp(n)^2\Bigr)\,D(\Om).
\end{equation}
\end{lemma}

\begin{proof}
Apply Lemma~\ref{lem:truncation} to get $\Otilde$ satisfying
\eqref{eq:trunc-vol}--\eqref{eq:trunc-A}; by Remark~\ref{rem:translate}
assume $\Otilde\subset B_{d(n)}\subset B_{R_*(n)}$. Theorem
\ref{thm:SV-bounded} at $R=R_*(n)$ gives
$E(\Otilde)-E(B_1)\ge c_{\mathrm{SV}}(n,R_*)\Fr(\Otilde)^2$; dividing by
$\omega_n^{(n+2)/n}$ and using $\abs{\Otilde}=\omega_n$,
\begin{equation}\label{eq:D-tilde-A}
   D(\Otilde)\;\ge\;\widetilde c_{\mathrm{SV}}(n)\,\Fr(\Otilde)^2 .
\end{equation}
%% JUSTIFICATION.  D(\Om)\le\delta(n) is the hypothesis of
%% Lemma~\ref{lem:truncation}, so the surrogate \Otilde exists with the
%% four properties.  After translation \Otilde\subset B_{d(n)}\subset
%% B_{R_*(n)} (as R_*\ge d).  Theorem~\ref{thm:SV-bounded} at R=R_*
%% applies (R_*\ge2, \Otilde\subset B_{R_*}, |\Otilde|=\omega_n):
%%   E(\Otilde)-E(B_1)\ge c_{\mathrm{SV}}(n,R_*)\Fr(\Otilde)^2.
%% Divide by \omega_n^{(n+2)/n}; the left becomes D(\Otilde) (unit
%% volume), the right \widetilde c_{\mathrm{SV}}(n)\Fr(\Otilde)^2.
Combining \eqref{eq:D-tilde-A} with \eqref{eq:trunc-D},
\[
   \Fr(\Otilde)^2\le\frac{D(\Otilde)}{\widetilde c_{\mathrm{SV}}}
   \le\frac{C_\sharp}{\widetilde c_{\mathrm{SV}}}\,D(\Om).
\]
%% JUSTIFICATION.  From \eqref{eq:D-tilde-A}, \Fr(\Otilde)^2\le
%% D(\Otilde)/\widetilde c_{\mathrm{SV}}; from \eqref{eq:trunc-D},
%% D(\Otilde)\le C_\sharp D(\Om).  Chain.
Now use \eqref{eq:trunc-A}, the elementary
$(a+b)^2\le2a^2+2b^2$, and $D(\Om)\le\delta(n)\le1$ (so
$D(\Om)^2\le D(\Om)$):
\begin{align*}
   \Fr(\Om)^2
   &\le\bigl(\Fr(\Otilde)+C_\sharp D(\Om)\bigr)^2
    \le2\Fr(\Otilde)^2+2C_\sharp^2D(\Om)^2\\
   &\le\frac{2C_\sharp}{\widetilde c_{\mathrm{SV}}}D(\Om)+2C_\sharp^2D(\Om)
    \le\Bigl(\frac{4C_\sharp}{\widetilde c_{\mathrm{SV}}}+4C_\sharp^2\Bigr)D(\Om),
\end{align*}
which is \eqref{eq:near-Fr}.
%% JUSTIFICATION, line by line.
%%  (1) \eqref{eq:trunc-A}: \Fr(\Om)\le\Fr(\Otilde)+C_\sharp D, square it.
%%  (2) (a+b)^2\le2a^2+2b^2 with a=\Fr(\Otilde), b=C_\sharp D:
%%      \le2\Fr(\Otilde)^2+2C_\sharp^2 D^2.
%%  (3) \Fr(\Otilde)^2\le(C_\sharp/\widetilde c_{\mathrm{SV}})D (above),
%%      and D^2\le D (since 0\le D\le1): \le2(C_\sharp/\widetilde
%%      c_{\mathrm{SV}})D+2C_\sharp^2 D.
%%  (4) factor D and bound 2\cdot\le4\cdot (just enlarge constants for a
%%      clean final form): \le(4C_\sharp/\widetilde c_{\mathrm{SV}}
%%      +4C_\sharp^2)D.  This is \eqref{eq:near-Fr}.
\end{proof}

\begin{corollary}[Small-deficit quadratic bound]\label{cor:near}
With
\begin{equation}\label{eq:cnear}
   c_{\mathrm{near}}(n):=
   \frac{\omega_n^{(n+2)/n}}
        {\,4C_\sharp(n)\bigl(1/\widetilde c_{\mathrm{SV}}(n)+C_\sharp(n)\bigr)}>0,
\end{equation}
every open $\Om$ with $\abs{\Om}=\omega_n$ and $D(\Om)\le\delta(n)$
satisfies $E(\Om)-E(B_1)\ge c_{\mathrm{near}}(n)\,\Fr(\Om)^2$.
\end{corollary}

\begin{proof}
Rearrange \eqref{eq:near-Fr} to
$D(\Om)\ge\Fr(\Om)^2/\bigl(4C_\sharp/\widetilde c_{\mathrm{SV}}
+4C_\sharp^2\bigr)$ and multiply by $\omega_n^{(n+2)/n}$, using
$E(\Om)-E(B_1)=\omega_n^{(n+2)/n}D(\Om)$.
%% JUSTIFICATION.  \eqref{eq:near-Fr} is \Fr^2\le A\cdot D with
%% A:=4C_\sharp/\widetilde c_{\mathrm{SV}}+4C_\sharp^2>0; invert to
%% D\ge\Fr^2/A.  Multiply by \omega_n^{(n+2)/n}: E(\Om)-E(B_1)=
%% \omega_n^{(n+2)/n}D\ge(\omega_n^{(n+2)/n}/A)\Fr^2.  Note
%% \omega_n^{(n+2)/n}/A=c_{\mathrm{near}}(n) after factoring 4C_\sharp out
%% of A: A=4C_\sharp(1/\widetilde c_{\mathrm{SV}}+C_\sharp).
\end{proof}

\subsection{Proof of the global theorem}\label{ssec:global-proof}

We now glue the far and near branches, then remove the volume
normalisation.

\begin{proof}[Proof of Theorem~\ref{thm:SV-global}]
Set
\begin{equation}\label{eq:csv-glob}
   c_{\mathrm{SV}}^{\mathrm{glob}}(n)
   :=\min\Bigl\{\,c_{\mathrm{near}}(n),\;
   \omega_n^{(n+2)/n}\,\tfrac{\delta(n)}{16}\,\Bigr\}>0 .
\end{equation}
Let $\Om$ be open with $\abs{\Om}=\omega_n$.
%% JUSTIFICATION.  Both arguments to the min are positive (c_{\mathrm{near}}
%% >0 from \eqref{eq:cnear}; \delta(n)>0), so c_{\mathrm{SV}}^{\mathrm{glob}}
%% (n)>0.  Taking the min lets a SINGLE constant work in both regimes.

\emph{Case A: $D(\Om)\ge\delta(n)/4$.} Lemma~\ref{lem:far} gives
$E(\Om)-E(B_1)\ge\omega_n^{(n+2)/n}(\delta(n)/16)\Fr(\Om)^2
\ge c_{\mathrm{SV}}^{\mathrm{glob}}(n)\Fr(\Om)^2$.
%% JUSTIFICATION.  Lemma~\ref{lem:far} applies (D\ge\delta/4) and gives
%% the first \ge; the second is because \omega_n^{(n+2)/n}\delta/16\ge
%% \min\{...\}=c_{\mathrm{SV}}^{\mathrm{glob}}.

\emph{Case B: $D(\Om)\le\delta(n)$} (which includes
$D(\Om)\le\delta(n)/4$). Corollary~\ref{cor:near} gives
$E(\Om)-E(B_1)\ge c_{\mathrm{near}}(n)\Fr(\Om)^2
\ge c_{\mathrm{SV}}^{\mathrm{glob}}(n)\Fr(\Om)^2$.
%% JUSTIFICATION.  Corollary~\ref{cor:near} applies (D\le\delta) and
%% gives the first \ge; the second because c_{\mathrm{near}}\ge
%% \min\{...\}=c_{\mathrm{SV}}^{\mathrm{glob}}.

The two cases overlap on $\delta(n)/4\le D(\Om)\le\delta(n)$ and together
cover $[0,\infty)$; in either case \eqref{eq:SV-global} holds with
$c_{\mathrm{SV}}^{\mathrm{glob}}(n)$ as in \eqref{eq:csv-glob}, a
dimensional constant because $\delta(n)$, $c_{\mathrm{near}}(n)$ depend
only on $n$ (recall $R_*(n)$ and hence $c_{\mathrm{SV}}(n,R_*(n))$ are
dimensional).
%% JUSTIFICATION (coverage).  Case A covers [\delta/4,\infty), Case B
%% covers [0,\delta]\supset[0,\delta/4]; their union is [0,\infty).  Any
%% D(\Om)\ge0 falls in at least one case, so \eqref{eq:SV-global} holds
%% for every unit-volume \Om.  All constants are functions of n alone.

\emph{General volume.} For arbitrary finite-measure $\Om$ put
$\Om_*:=(\omega_n/\abs{\Om})^{1/n}\Om$, so $\abs{\Om_*}=\omega_n$,
$\Fr(\Om_*)=\Fr(\Om)$ and $D(\Om_*)=D(\Om)$ by scale invariance of
$\Fr$ and $D$. Applying \eqref{eq:SV-global} to $\Om_*$ reads
$D(\Om)\ge\omega_n^{-(n+2)/n}c_{\mathrm{SV}}^{\mathrm{glob}}(n)\Fr(\Om)^2$.
%% JUSTIFICATION (the rescaling).  With t:=(\omega_n/|\Om|)^{1/n},
%% |\Om_*|=t^n|\Om|=\omega_n.  \Fr and D are scale invariant (dilation
%% does not change normalized asymmetry, nor the scale-invariant deficit),
%% so \Fr(\Om_*)=\Fr(\Om), D(\Om_*)=D(\Om).  \eqref{eq:SV-global} for the
%% unit-volume set \Om_*: E(\Om_*)-E(B_1)\ge c^{\mathrm{glob}}\Fr(\Om_*)^2;
%% divide by \omega_n^{(n+2)/n} (=\omega_n^{(n+2)/n}, unit volume), using
%% D(\Om_*)=\omega_n^{-(n+2)/n}(E(\Om_*)-E(B_1)):
%%   D(\Om_*)\ge\omega_n^{-(n+2)/n}c^{\mathrm{glob}}\Fr(\Om_*)^2,
%% i.e., by invariance, D(\Om)\ge\omega_n^{-(n+2)/n}c^{\mathrm{glob}}\Fr(\Om)^2.
Multiplying by $\abs{\Om}^{(n+2)/n}$ and using
$E(\Om)\abs{\Om}^{-(n+2)/n}-E(B^*)\abs{B^*}^{-(n+2)/n}=D(\Om)$ with
$\abs{B^*}=\abs{\Om}$, then $E=-T/2$, gives
\[
   T(B^*)-T(\Om)=2\bigl(E(\Om)-E(B^*)\bigr)
   \ge2\,c_{\mathrm{SV}}^{\mathrm{glob}}(n)
   \Bigl(\tfrac{\abs{\Om}}{\omega_n}\Bigr)^{(n+2)/n}\Fr(\Om)^2,
\]
which is \eqref{eq:T-global}.
%% JUSTIFICATION (undoing the normalisation).  Multiply
%% D(\Om)\ge\omega_n^{-(n+2)/n}c^{\mathrm{glob}}\Fr(\Om)^2 by
%% |\Om|^{(n+2)/n}:
%%   |\Om|^{(n+2)/n}D(\Om)\ge(|\Om|/\omega_n)^{(n+2)/n}c^{\mathrm{glob}}\Fr^2.
%% By \eqref{eq:deficit-def} with |B^*|=|\Om|,
%%   |\Om|^{(n+2)/n}D(\Om)=E(\Om)-E(B^*)(|\Om|/|B^*|)^{(n+2)/n}=E(\Om)-E(B^*),
%% wait: D(\Om)=E(\Om)|\Om|^{-(n+2)/n}-E(B^*)|B^*|^{-(n+2)/n} and
%% |B^*|=|\Om|, so |\Om|^{(n+2)/n}D(\Om)=E(\Om)-E(B^*).  Hence
%%   E(\Om)-E(B^*)\ge c^{\mathrm{glob}}(|\Om|/\omega_n)^{(n+2)/n}\Fr^2.
%% Finally E=-T/2 (\eqref{eq:energy-identities}) gives E(\Om)-E(B^*)
%% =\tfrac12(T(B^*)-T(\Om)); multiply by 2:
%%   T(B^*)-T(\Om)=2(E(\Om)-E(B^*))\ge2c^{\mathrm{glob}}
%%     (|\Om|/\omega_n)^{(n+2)/n}\Fr^2.  This is \eqref{eq:T-global}.
\end{proof}

\begin{remark}[Sharpness of the exponent]
The exponent $2$ on $\Fr$ in \eqref{eq:SV-global} is optimal: it is
saturated by smooth, volume-preserving ellipsoidal perturbations of
$B_1$, for which both $E(\Om)-E(B_1)$ and $\Fr(\Om)^2$ are of the same
quadratic order in the perturbation parameter.
%% JUSTIFICATION.  For an ellipsoidal perturbation \Om_t of B_1 with
%% volume preserved and perturbation parameter t\to0, a second-order shape
%% calculation gives E(\Om_t)-E(B_1)\sim c_1 t^2 (the first variation
%% vanishes at the ball by Saint--Venant optimality, the second variation
%% is a nonzero quadratic form on the nonradial modes) and \Fr(\Om_t)\sim
%% c_2 t (the asymmetry is linear in the displacement), so \Fr(\Om_t)^2\sim
%% c_2^2 t^2.  Thus the ratio (E(\Om_t)-E(B_1))/\Fr(\Om_t)^2 stays bounded
%% as t\to0; no exponent larger than 2 can hold.  (This shows optimality
%% of the EXPONENT; it does not bear on the value of the constant.)
\end{remark}

%=====================================================================
%  END of exhaustive bounded-reduction section.
%=====================================================================
